\documentclass[11pt]{article}
\usepackage[T1]{fontenc}
\usepackage[utf8]{inputenc}
\usepackage{lmodern,microtype}
\usepackage[margin=0.85in]{geometry}
\usepackage{amsmath,amssymb,amsthm,mathtools,mathrsfs}
\usepackage{array,enumitem,xcolor}
\usepackage{hyperref}
\usepackage{aliascnt}
\usepackage[nameinlink,capitalise,noabbrev]{cleveref}
\usepackage{fancyhdr}
\hypersetup{colorlinks=true,linkcolor=blue!40!black,citecolor=blue!40!black,
  urlcolor=blue!40!black,pdftitle={Standalone Lean 4 blueprint: lines 2183 and 2193},
  pdfauthor={}}
\setlength{\parindent}{0pt}
\setlength{\parskip}{0.45em}
\setlength{\headheight}{14pt}
\setlength{\emergencystretch}{2em}
\pagestyle{fancy}
\fancyhf{}
\lhead{Standalone Lean 4 blueprint}
\rhead{Lines 2183 and 2193}
\cfoot{\thepage}
\newtheorem{theorem}{Theorem}[section]
\newaliascnt{proposition}{theorem}
\newtheorem{proposition}[proposition]{Proposition}
\aliascntresetthe{proposition}
\newaliascnt{lemma}{theorem}
\newtheorem{lemma}[lemma]{Lemma}
\aliascntresetthe{lemma}
\newaliascnt{corollary}{theorem}
\newtheorem{corollary}[corollary]{Corollary}
\aliascntresetthe{corollary}
\theoremstyle{definition}
\newaliascnt{definition}{theorem}
\newtheorem{definition}[definition]{Definition}
\aliascntresetthe{definition}
\newaliascnt{remark}{theorem}
\newtheorem{remark}[remark]{Remark}
\aliascntresetthe{remark}
\newcommand{\ee}[1]{\mathrm e\!\left(#1\right)}
\newcommand{\pid}[1]{\nolinkurl{#1}}
\begin{document}
\begin{center}
{\LARGE\bfseries Fej\'er--Gowers cube patch}\\[7pt]
{\large A standalone Lean 4 mathematical blueprint}\\[5pt]
{\normalsize Original lines 2183 and 2193 of \texttt{lem:fejer-vdc}}\\[8pt]
{\small The real monomial system $t,t^2,t^3$: proof and necessary dependency closure}
\end{center}
\vspace{0.4em}
\setcounter{tocdepth}{2}
\tableofcontents
\clearpage
\section*{Scope, mathematical changes, and proof status}
\addcontentsline{toc}{section}{Scope and replacement map}
This standalone mathematical blueprint addresses only the two blocked passages at original
lines 2183 and 2193, inside \pid{lem:fejer-vdc} in
\pid{task_2_smoothingineq3d_blueprint_updated.tex}.
The completed repairs at lines 2084 and 2086 are outside its scope and must remain unchanged.
No subunit-endpoint proof or replacement of those earlier passages is included.
Line numbers identify the original upload; use theorem labels to locate the passages in a
manuscript that already contains earlier edits.

\paragraph{Line 2183: the unproved assertion.}
The blocked sentence says that repeated van der Corput and Cauchy--Schwarz bound
``every polynomial average obtained from the monomials $t,t^2,t^3$'' by a finite power of the
local uniformity norm of a selected input, with a boundary error.
It does not specify a polynomial state or prove that the selected input survives in the
configuration needed by the terminal estimate.

\paragraph{Line 2193: the missing removal and cube induction.}
The blocked paragraph says that ``Cauchy--Schwarz removes'' each unselected factor and that
an induction ``gives exactly the cube product.''
It does not identify the block removed at each step, give the induction invariant, or derive
the final configuration to which \pid{cfgHeadInt_cube} applies.

\begin{center}
\begin{tabular}{>{\raggedright\arraybackslash}p{.20\linewidth}>{\raggedright\arraybackslash}p{.72\linewidth}}
\hline
Primary blocker & Replacement \\
\hline
2183: unrestricted assertion & Retain the scalar van der Corput lemma.  Replace the
unqualified assertion about all generated configurations by the actual monomial conclusion,
proved through an independent real inverse module.\\
2193: asserted cube endpoint & Explicit coordinate-doubling states, a protected-factor
polynomial invariant, and, for the public every-input conclusion, an explicitly expanded
order-two cube.\\
\hline
\end{tabular}
\end{center}

\paragraph{Necessary dependent changes, not additional blockers.}
The repair includes the affected \pid{lem:pet-reduction} and the inverse-theorem proof because
they use the missing implication.  The every-input local-uniformity result has $s_*=2$;
the public inverse theorem is rebuilt from the same independently proved positive-energy
theorem, rather than used circularly to supply its own premise.
The old unrestricted \pid{lem:degree-lowering-zero} is removed from this proof path.
These supporting arguments are included here, so the document does not require an earlier
terminal patch or a separate unproved input-producing certificate.

\paragraph{What changes in mathematical scope.}
The last assertion of the old Fej\'er lemma is not retained as a theorem about arbitrary
generated configurations.  Its needed consequence for $(t,t^2,t^3)$ is proved below, with
the original every-$j$ quantifier and physical-scale bounds.  This is an explicit statement
repair, not a proof of the old, more general wording.

\paragraph{Sources versus added derivations.}
The source is the supplied version of Kosz--Mirek--Peluse--Wan--Wright \cite{KoszEtAl}.
The real inverse module follows its organization: highest-degree uniformity control;
conditional degree lowering and the major-arc induction; lowest-degree inverse information;
then the general-index reduction.  The protected-leading-coefficient PET invariant, the
quantitative coordinate-doubling lemmas, and the final low-frequency-to-$U^2$ bridge are
added derivations in this blueprint.  They are not attributed verbatim to the paper.
In particular, we do not invoke the integer statement of its reference [63] as a theorem over
$\mathbb R$, and do not assume the paper's Theorem 4.30 in order to prove our own inverse theorem.

\paragraph{Lean boundary.}
This is a proposed mathematical proof blueprint, not a compiler-checked Lean library.
Project declaration names below are proposed names unless the user has already supplied them.
The existing \pid{fejer_vdc'}, \pid{locUnifPow}, and \pid{cfgHeadInt_cube} are retained.
Their precise Lean binders were not supplied; the document specifies mathematical inputs
and proof obligations, not invented API signatures.  No inverse theorem, PET certificate,
or uniformity estimate is added as a hypothesis to a public conclusion.

\paragraph{Background library.}
The analytic background is Lebesgue integration, Fubini--Tonelli, Cauchy--Schwarz and H\"older,
standard convergence theorems, translation and dilation changes of variables,
elementary polynomial algebra and calculus, and Plancherel.  These are library prerequisites,
not new analytic assumptions specific to this problem.  The nonstandard consequences used
in the repair are developed below.

\section{Conventions and quantitative bookkeeping}\label{patch:conventions}
We use $\ee u=\exp(2\pi i u)$, $\langle f,g\rangle=\int f\overline g$, and
$\widehat f(\xi)=\int f(x)\ee{-x\cdot\xi}\,dx$.  Write $\mathcal C$ for complex conjugation.
For $N>0$,
\[
 A_N(f_1,f_2,f_3)(x)=\frac1N\int_0^N\prod_{i=1}^3f_i(x+t^ie_i)\,dt.
\]
and $\Lambda_N=\langle A_N(f_1,f_2,f_3),f_0\rangle$.
The notation $\|f\|_p$ denotes the Lebesgue $L^p$ norm.
Fix an even real smooth cutoff $\eta$ with
$\mathbf1_{[-1/4,1/4]}\leq\eta\leq\mathbf1_{(-1/2,1/2)}$.
The operator $P_L^{(j)}$ has multiplier $\eta(\xi_j/L)$.

Work first with Borel representatives, pointwise bounded by the stated bounds and zero off
the stated boxes.  Almost-everywhere versions follow at the end: for a null set $E$ and any
measurable translation $v(t)$,
\[
 \int\!\int\mathbf1_E(x+v(t))\,dx\,dt=0
\]
by translation invariance and Tonelli.  The same argument handles any finite number of shift
parameters.  Thus changing representatives does not change any correlation or cube integral.
All auxiliary phase functions can be chosen Borel after discarding the standard Fubini null
sets.

\subsection{A precise meaning for the power losses}
For an integer $c\geq1$, set
\[
 \ell_c(\delta)=c^{-1}\delta^c,\qquad u_c(\delta)=c\delta^{-c}
 \qquad(0<\delta\leq1).
\]
An assertion below of the form $X\geq\delta^{O(1)}S$ means that the proof returns an integer
$c$, depending only on the fixed degrees, support/coefficient budgets, and any fixed cube order,
with $X\geq\ell_c(\delta)S$.  Similarly $X\leq\delta^{-O(1)}S$ returns $c$ with
$X\leq u_c(\delta)S$.  It never means that the exponent may depend on a function, a selected
section, $N$, or the numerical value of $\delta$.

For implementation, retain the returned integers rather than using asymptotic notation.
The following rules suffice to compose every budget in the proof.
\begin{lemma}[Budget arithmetic and popularity]\label{patch:budgets}
Finite sums, products, quotients with a positive lower bound, and fixed natural powers of
power-controlled positive quantities are power-controlled.  A fixed positive root of a
power lower bound is again a power lower bound.  If a finite measure space has measure at most
$M$, $0\leq F\leq B$, and $\int F\geq a>0$, then
\[
 E=\{F\geq a/(2M)\}\quad\hbox{satisfies}\quad \mu(E)\geq a/(2B).
\]
Removing a set of measure at most $a/(2B)$ loses at most half the lower bound.  A partition
into $K$ measurable cells has a cell carrying at least $1/K$ of a nonnegative total mass.
\end{lemma}
\begin{proof}
For upper bounds choose, for example, an integer at least both the new coefficient and the
new exponent.  For a product of $u_a$ and $u_b$ it suffices to take
$c\geq\max(ab,a+b)$; for a product of $\ell_a$ and $\ell_b$ the same choice works.
For $r$-th powers take $c\geq\max(a^r,ra)$.  Bound $X/Y$ using an upper bound for $X$
and a positive lower bound for $Y$, or the reverse for a lower bound.  Round all constants
upward by the Archimedean property.  Since $0<\delta\leq1$, increasing a lower-bound exponent
weakens it and increasing an upper-bound exponent enlarges it.  These observations also
handle roots and rational fixed exponents.
For popularity, the integral on $E^c$ is at most $a/2$, whereas the integral on $E$ is at most
$B\mu(E)$.  The remaining assertions follow from integral monotonicity and a finite sum.
\end{proof}

Whenever a proof chooses a finer parameter, it uses these explicit recipes:
\[
 \text{discarded measure}\leq\frac{\text{retained mass}}{4\,\text{integrand bound}},
 \qquad
 \text{phase error}\leq\frac{\text{retained mass}}{4\,\text{integration-volume bound}}.
\]
For a phase perturbation use $|\ee u-\ee v|\leq2\pi|u-v|$.
A bounded interval of frequencies of length $2M$ can be partitioned into at most
$1+\lceil2M/\tau\rceil$ cells of width $\tau$.  All such counts below are power-controlled.
The iteration counts are bounded independently of $\delta$ before these recipes are applied.

\paragraph{Scale convention in the inverse module.}
For $1\leq d_1<\cdots<d_k\leq3$, write $D_r=\sum_{i\leq r}d_i$, $D_0=0$ and $D=D_k$.
A budgeted support box has sides $[-u_c(\delta)N^{d_i},u_c(\delta)N^{d_i}]$.
A polynomial $P_i$ is admissible with budget $c$ when
\[
 \deg P_i=d_i,\qquad
 \ell_c(\delta)\leq|\operatorname{lc}P_i|\leq u_c(\delta),\qquad
 |[t^r]P_i|\leq u_c(\delta)N^{d_i-r}\quad(0\leq r<d_i).
\]
These are the paper's admissibility and enlarged-support hypotheses, with a single integer
budget enlarged as needed.  Setting $t=Nu$ and $x_i=N^{d_i}y_i$ reduces all coefficient,
volume, and sublevel estimates to bounded parameter intervals with power-controlled
constants.  This change has determinant $N^D$; the powers of $N$ in subsequent statements
are therefore exact, not part of the budget.


\section{An independent real inverse module}\label{sec:internal-kosz-adjoint}
The source organization is \cite[\S\S4.2--4.5]{KoszEtAl}.  The public every-input uniformity
lemma is \emph{not} a premise of this section.  The only polynomial uniformity estimate used
inside the section is the highest-active-degree estimate proved below, including its
parameter-shifted version.  We work with all increasing degree lists contained in $\{1,2,3\}$,
not just the list $(1,2,3)$: shorter lists and lower coefficients are needed by the slicing argument.

\subsection{Fej\'er measures and the existing cube quantity}
\begin{definition}[Fej\'er differences and local uniformity]\label{def:local-uniformity}
For $H>0$ let $\kappa_H(h)=H^{-1}(1-|h|/H)_+$.
For $v\in\mathbb R^3$ use the manuscript's difference
$\Delta_{h;v}f(x)=f(x+hv)\overline{f(x)}$.
The existing quantity is
\[
 \mathcal U^s_{j,H}(f)^{2^s}
 =N^{-6}\int_{\mathbb R^3}\int_{\mathbb R^s}
       \Delta_{h_1;e_j}\cdots\Delta_{h_s;e_j}f(x)
           \prod_i\kappa_H(h_i)\,d\mathbf h\,dx.
\]
In particular this definition is not replaced by a global Gowers norm or by Fourier energy.
\end{definition}
For calculations let
\[
 C_s(f;x,\mathbf h)=\prod_{\omega\in\{0,1\}^s}
       \mathcal C^{|\omega|}f\bigl(x+(\omega\cdot\mathbf h)e_j\bigr),
 \qquad C_0(f;x)=f(x).
\]
Then
\begin{equation}\label{patch:cube-succ}
 C_{s+1}(f;x,\mathbf h,h)=C_s(f;x,\mathbf h)
                                  \overline{C_s(f;x+he_j,\mathbf h)}.
\end{equation}
For the manuscript's difference convention the displayed parity cube is globally conjugated
when $s$ is odd.  Its integrated Fej\'er power is real, so the two conventions give the same
quantity.  This is a convention bridge only; retain the definition and convention used by
\pid{locUnifPow} and \pid{cfgHeadInt_cube}.

In ambient dimension $n$, with an explicit positive spatial normalizer $V$, write
\[
 Q_{s,H,V}^{j}(f)=V^{-1}\operatorname{Re}\int\!\int C_s(f;x,\mathbf h)
                                   \prod_i\kappa_H(h_i)\,d\mathbf h\,dx.
\]
The notation with unequal radii, $Q_{\mathbf L,V}^{j}$, means the same expression with
$\prod_i\kappa_{L_i}(h_i)$.  These are proof-local normalizations.  The public order-$s$
quantity is exactly $Q_{s,H,N^6}^{j}$, not $Q_{s,H,|B_N|}^{j}$.
Define $T_{H,j}u(x)=H^{-1}\int_0^H u(x+te_j)\,dt$.

\begin{lemma}[Squares, positivity, and appending a coordinate]\label{patch:fejer-squares}
The difference of two independent uniform variables in $[0,H]$ has density $\kappa_H$.
Consequently
\[
 \int\!\int u(x)\overline{u(x+he_j)}\kappa_H(h)\,dh\,dx
       =\|T_{H,j}u\|_2^2.
\]
For $s\geq1$, every order-$s$ cube power above is real and nonnegative.  If $|f|\leq1$ and $f$ is supported in a
rectangular box of volume $V$, then $0\leq Q_{s,H,V}^{j}(f)\leq1$.
If the box has $j$-th side length $S$, then
\begin{equation}\label{patch:raise-order}
 \bigl(Q_{s,H,V}^{j}(f)\bigr)^2
       \leq(1+H/S)Q_{s+1,H,V}^{j}(f).
\end{equation}
\end{lemma}
\begin{proof}
Expand the square of $T_{H,j}u$, use translation in $x$, and apply Fubini.  The density of
$a-b$ for $a,b\in[0,H]$ is $H^{-2}|[0,H]\cap([0,H]-h)|=\kappa_H(h)$.
Apply this identity to $u=C_{s-1}(f;\cdot,\mathbf h)$ for the last difference.  It proves
reality and positivity before any comparison of kernels is made.  The zero vertex gives
$|C_s(f;x,\mathbf h)|\leq\mathbf1_B(x)$, so the upper bound is one.
For \eqref{patch:raise-order}, $u=C_s$ is supported in $B$, its full integral is unchanged
by $T_{H,j}$, and $T_{H,j}u$ is supported in a box of volume $V(1+H/S)$.
Thus $|\int u|^2\leq V(1+H/S)\|T_{H,j}u\|_2^2$.  Average over the first $s$ increments,
apply Jensen, and use \eqref{patch:cube-succ}.  The same argument applies to translated
support boxes.  Translation and conjugation of $f$ preserve every integrated cube power.
\end{proof}

\begin{lemma}[Fej\'er van der Corput]\label{lem:fejer-vdc}
Let $I$ have length $N>0$, $0<H\leq N/4$, and $F$ be measurable and $1$-bounded.  Then
\[
 \left|\frac1N\int_I F(t)\,dt\right|^2
 \leq4\int\kappa_H(h)
       \left|\frac1N\int_{I\cap(I-h)}F(t+h)\overline{F(t)}\,dt\right|\,dh
       +\frac{8H}{N}.
\]
\end{lemma}
\begin{proof}
Put $W=\mathbf1_IF$ and $S(u)=H^{-1}\int_0^H W(u+a)\,da$.
Then $\int S=\int_I F$ and $S$ is supported in an interval of length $N+H$.
Cauchy--Schwarz followed by square expansion gives the stronger upper bound
$(1+H/N)\int\kappa_H(h)|N^{-1}\int_{I\cap(I-h)}F(t+h)\overline{F(t)}\,dt|\,dh$.
Use $1+H/N\leq2$.  The displayed version is therefore valid; it is the existing
\pid{fejer_vdc'} interface.  The plus sign before $8H/N$ is explicit.
\end{proof}
No unrestricted polynomial-to-uniformity assertion is included in this scalar lemma.

\begin{lemma}[Signed integrated removal]\label{patch:signed-vdc}
Let $I$ have length $N>0$ and $0<H\leq N/4$.  Suppose $V>0$, $\int|g|^2\leq V$, and $\int|F(x,t)|^2\,dx\leq V$ for the relevant
parameter values.  Assume joint measurability and the integrability supplied, in applications,
by bounded factors and their finite-support envelopes.  Set
$A=V^{-1}\int g(x)\mathbb E_{t\in I}F(x,t)\,dx$.  Then
\begin{equation}\label{patch:signed-vdc-eq}
 |A|^2\leq2\operatorname{Re}\int\kappa_H(h)
      \frac1V\int\mathbb E_{t\in I}F(x,t)\overline{F(x,t+h)}\,dx\,dh
      +2H/N.
\end{equation}
\end{lemma}
\begin{proof}
First remove exactly $g$, which is independent of $t$, by Cauchy--Schwarz in $x$:
$|A|^2\leq V^{-1}\int|\mathbb E_I F(x,t)|^2dx$.
Apply the window construction in the preceding proof with $x$ retained and integrate in $x$.
It gives $(1+H/N)$ times the real part of the averaged cut-off correlation.  This whole
averaged correlation is nonnegative because it came from an integrated square.
Its replacement of $I\cap(I-h)$ by $I$ costs at most $|h|/N$: indeed
$\int|F(x,t)\overline{F(x,t+h)}|dx\leq V$ by Cauchy--Schwarz.
Bound $1+H/N$ by two and average $|h|\leq H$.
This proves \eqref{patch:signed-vdc-eq}.  In particular, it is not obtained by deleting an
absolute value from \pid{fejer_vdc'}.
\end{proof}

\subsection{The elementary coordinate-removal induction}
\begin{lemma}[Cylinder Cauchy--Schwarz, with its state]\label{patch:cylinder}
Let $z$ be an auxiliary probability-space variable and $b_i$ independent probability-space
variables, $1\leq i\leq r$.  Let $v(z,\mathbf b)$ and $u_i(z,\mathbf b)$ be jointly measurable and $1$-bounded,
with $u_i$ independent of $b_i$.  Then
\[
 \left|\mathbb E_{z,\mathbf b}v(z,\mathbf b)\prod_{i=1}^r u_i(z,\mathbf b)\right|^{2^r}
 \leq\mathbb E_{z,\mathbf b^0,\mathbf b^1}
       \prod_{\omega\in\{0,1\}^r}\mathcal C^{|\omega|}v(z,\mathbf b^\omega).
\]
The right-hand side is real and nonnegative.  The same statement for finite nonprobability
spaces follows by writing every normalization factor explicitly.
\end{lemma}
\begin{proof}
After $a$ steps use independent pairs $b_i^0,b_i^1$ for $i\leq a$ and one variable $b_i$
for $i>a$.  Define
\[
 v_a=\prod_{\omega\in\{0,1\}^a}\mathcal C^{|\omega|}v(z,\mathbf b^\omega),
 \qquad
 u_{i,a}=\prod_{\omega\in\{0,1\}^a}\mathcal C^{|\omega|}u_i(z,\mathbf b^\omega)
 \quad(i>a),
\]
and $J_a=\mathbb E[v_a\prod_{i>a}u_{i,a}]$.
At step $a+1$ the removed block is exactly $u_{a+1,a}$.
It is independent of the only variable being doubled, $b_{a+1}$.
Write the average as an outer average of this block times the inner $b_{a+1}$ average.
Cauchy--Schwarz in all the outer variables and expansion of the inner square yield
$|J_a|^2\leq J_{a+1}$.  The new $J_{a+1}$ is nonnegative because it is that square average.
The finite-product equivalence $\{0,1\}^{a+1}\simeq\{0,1\}^a\times\{0,1\}$ proves
that the expanded factors are precisely the displayed next state; the new bit toggles
conjugation.  Induct on $a=0,\ldots,r-1$.  The auxiliary variable $z$ is never doubled.
\end{proof}
This is the reusable, fully specified Gowers--Cauchy--Schwarz induction used below.
It does not claim that a polynomial family is already a cube.

\subsection{Scalar polynomial oscillation and coefficient separation}
\begin{lemma}[Real polynomial oscillation]\label{lem:real-polynomial-oscillation}
For $1\leq d\leq3$ there is $C_d\geq1$ such that
\[
 \left|N^{-1}\int_0^N\ee{\sum_{r=1}^d a_rt^r}\,dt\right|
 \leq C_d\min\{1,(\max_r|a_r|N^r)^{-1/d}\}.
\]
When $\max_r|a_r|N^r=0$, interpret the minimum as one; in Lean use a separate zero case,
not a negative real power of zero.  Hence a lower bound $0<\varepsilon\leq1$ for this average bounds every nonconstant coefficient
by $|a_r|N^r\leq (C_d/\varepsilon)^d$.
\end{lemma}
\begin{proof}
Scale to $N=1$.  For a real polynomial $p$ of degree at most $d-1$ with a coefficient of
modulus at least $A$, a measurable sublevel set $E\subseteq[0,1]$ on which $|p|\leq\theta$
has length at most $C_d(\theta/A)^{1/(d-1)}$ when $d\geq2$.
To see this, choose $d$ points in $E$ with mutual separations at least $|E|/(2d)$ by successively
removing intervals about the previously chosen points (or use $d$ separated quantile blocks).
Lagrange interpolation then bounds every coefficient by $C_d\theta|E|^{-(d-1)}$.
Null sets and degree-zero polynomials are handled separately.
Apply this to $p=P'$ and $A=\max|a_r|$.  Split $[0,1]$ at the finitely many roots of
$P''$, $P'-\theta$, and $P'+\theta$.  On each remaining interval where $|P'|\geq\theta$,
$P'$ is monotone; integration by parts in $\ee{P}/(2\pi iP')$ bounds the integral by
$C_d/\theta$.  The total number of intervals is bounded in terms of $d$.
Choose $\theta=A^{1/d}$ when $A\geq1$.  The sublevel contribution and the integration-by-parts
contribution are both $O_d(A^{-1/d})$.  For $d=1$ integrate the exponential directly;
for $A<1$ use the bound one.  Scale back.
\end{proof}

\begin{lemma}[Distinct degrees give a triangular frequency bound]\label{patch:triangular}
For admissible $P_1,\ldots,P_k$, if
$|\mathbb E_{t\in[0,N]}\ee{\sum_i\xi_iP_i(t)}|\geq\delta^{O(1)}$, then
$\sum_iN^{d_i}|\xi_i|\leq\delta^{-O(1)}$.
\end{lemma}
\begin{proof}
Write $a_r=[t^r]\sum_i\xi_iP_i(t)$ and $z_i=N^{d_i}|\xi_i|$.
The preceding lemma bounds every $N^r|a_r|$ by a power of $\delta^{-1}$.
Starting at $i=k$ and descending, use the exact identity
\[
 \xi_i\operatorname{lc}(P_i)=a_{d_i}-\sum_{u>i}\xi_u[t^{d_i}]P_u.
\]
Consequently $z_i\leq u_c(\delta)\{M+u_c(\delta)\sum_{u>i}z_u\}$ for a known power bound $M$.
There are at most three steps.  Apply \cref{patch:budgets} to this explicit recursion.
\end{proof}


\subsection{The protected highest-degree PET producer}\label{patch:pet-producer}
This subsection supplies, rather than assumes, the polynomial input producer needed inside the
real inverse proof.  It is deliberately a highest-\emph{active}-degree result.  General selected
indices will be handled later by the independent slicing argument.

\subsubsection{Symbolic state and a leading-coefficient invariant}
Use $t$ for the averaging variable and $\mathbf h=(h_1,\ldots,h_r)$ for formal indeterminates.
Polynomials are elements of $(\mathbb R[\mathbf h])^n[t]$; degrees in this subsection are degrees
in $t$, \emph{before} evaluating the $h_i$.  A normal state consists of a finite ordered list
\[
 Q_1,\ldots,Q_L,\qquad Q_i(0,\mathbf h)=0,
\]
of nonzero distinct vector polynomials, with $Q_1$ of maximum degree, together with a bounded
function block $g_i$ for each $Q_i$ and a spatial block $g_0$.
A block is a finite list of translated/conjugated input factors.  Multiplicities are retained.
The distinguished block $g_1$ is required to contain \emph{one} translated or conjugated copy
of the protected input $f$, not a product containing $f$.

Fix the protected input's original degree $D\in\{2,3\}$, leading scalar coefficient $a\ne0$,
and coordinate $e_m$.  Retain the following invariant for $Q_1$ and for every nonzero
$Q_1-Q_i$, including $Q_0=0$:
\begin{equation}\label{patch:lead-invariant}
 \operatorname{lc}_t R=a\,p_R(\mathbf h)e_m,
 \quad p_R\in\mathbb Z[\mathbf h]\setminus\{0\},
 \quad p_R\text{ multilinear and homogeneous of degree }D-\deg_t R.
\end{equation}
Only the leading coefficients have this form.  The full polynomials may contain lower
coefficients in any active coordinate and with arbitrary admissible real values.

Initially the input polynomials are $P_i(t)e_i$ with increasing degrees and protected index $m$.
Then \eqref{patch:lead-invariant} holds with $p_R=1$ for the head and its differences with
lower-degree inputs.  It also holds for the shifted family
\begin{equation}\label{patch:initial-shift-family}
 P_i(t+\omega\cdot\mathbf h)e_i-P_i(\omega\cdot\mathbf h)e_i,
       \qquad \omega\in\{0,1\}^{r_0},\quad r_0\leq4,
\end{equation}
when the protected head is the highest input with $\omega=(1,\ldots,1)$.
For two shifts of that input, their difference has degree $D-1$ and leading coefficient
$aD[(\mathbf1-\omega)\cdot\mathbf h]e_m$; differences with lower inputs still have degree $D$.
Repeated normalized polynomials among nuisance blocks are grouped.  Since $D\geq2$, the
protected head has no duplicate in this initial family.

\begin{lemma}[A labelled update preserves the protected head]\label{patch:pet-update}
Suppose the current head degree is $d\geq2$.  For a pivot $p\in\{1,\ldots,L\}$ introduce one
fresh indeterminate $u$ and form the ordered children
\[
 Q_i(t+u)-Q_p(t),\quad Q_i(t)-Q_p(t)\qquad(1\leq i\leq L),
\]
putting the shifted head first.  Subtract each child's constant term in $t$, absorb that
constant as a translation of its function, collect zero polynomials in the new spatial block,
and multiply the function lists attached to identical remaining polynomials.
If the shifted head has maximum degree, the new state is normal, its head block is a
singleton of origin $f$, and \eqref{patch:lead-invariant} is preserved.
\end{lemma}
\begin{proof}
Let $\delta_u Q_1=Q_1(t+u)-Q_1(t)$.  Its degree is $d-1$ and leading coefficient
$d u\operatorname{lc}_tQ_1$.  For every $i$, compare the new head with an unshifted child:
\[
 [Q_1(t+u)-Q_p(t)]-[Q_i(t)-Q_p(t)]
       =\delta_uQ_1+(Q_1-Q_i).
\]
If the degrees of the two summands differ, take the larger.  If they agree, the leading
coefficient is the sum of an old polynomial in $\mathbf h$ and $du$ times another nonzero
old polynomial.  It cannot vanish because $u$ is a fresh indeterminate.  The sum is still
$a$ times a nonzero integer multilinear polynomial with the asserted homogeneous degree.
For a shifted child the difference is simply $(Q_1-Q_i)(t+u)$, whose leading coefficient is
unchanged.  Apply the first calculation with $i=p$ to the new head itself.
Subtracting constants does not change these nonconstant leading terms.
In particular the head differs nontrivially from every other child, and is nonconstant even
when $p=1$.  Therefore grouping never multiplies another factor into the protected head.
The shifted child conjugates its function, whereas an unshifted child retains its conjugation
flag.  Translation by its constant term preserves origin and singleton status.
\end{proof}
If all integer polynomials in the invariant have coefficient $\ell^1$-norm at most $M$,
the proof gives the bound $(D+1)M$ after an update.  Initially one can take
$M=1+2Dr_0$.  These are numerical bounds independent of the values of the real lower coefficients.

\subsubsection{The pivot is specified, and termination has a uniform bound}
Let $w_q$ count distinct leading coefficients of degree-$q$ polynomials in the normal state.
Order $(w_D,\ldots,w_1)$ lexicographically.  Stop when every degree is at most one.
Otherwise let $d=\deg_tQ_1$ and let $l$ be the minimum degree present.  Choose the first
index in the ordered list satisfying the relevant case:
\[
 \begin{array}{ll}
 l<d:&\deg Q_p=l;\\
 l=d,\ w_d>1:&\operatorname{lc}Q_p\ne\operatorname{lc}Q_1;\\
 l=d,\ w_d=1:&p=1.
 \end{array}
\]
For the first two cases the head retains maximum degree; for the third it has degree $d-1$
and every child has degree at most $d-1$.  For the selected degree $l$, the new leading
classes are the nonzero differences with the pivot leading coefficient.  Thus
\[
 w'_q=w_q\ (q>l),\qquad w'_l=w_l-1,\qquad 0\leq w'_q\leq2L\ (q<l).
\]
The type strictly decreases and the new length is at most $2L$.
The argument is in the formal polynomial ring; numerical exceptional parameter values do
not change the choice of pivot.

For an actual uniform iteration bound define $B(w,L)$ by well-founded recursion on the type,
with its value a function of the length budget $L$:
\[
 B(w,L)=0\quad\hbox{at a linear type},\qquad
 B(w,L)=1+\max_{w'\in\mathscr C(w,L)}B(w',2L)\quad\hbox{otherwise},
\]
where $\mathscr C(w,L)$ is the finite set of types with the preceding displayed bounds for
the chosen $l$.  Each $w'<w$, so this is a well-founded definition.  The usual induction
proves that it bounds the execution on every family of length at most $L$ and type $w$.
If the initial length is at most $L_0\leq3\cdot2^4$, take the finite maximum of $B(w,L_0)$
over $0\leq w_q\leq L_0$.  Call it $R_0$.  The terminal length is at most $L_1=2^{R_0}L_0$.
This proves uniform complexity, not merely termination of each individual execution.

\paragraph{Full coefficient and translation bounds.}
For a fixed spatial coordinate $i$, retain a bound $M N^{d_i-q}$ for the coefficient of
$t^q$ of every current polynomial and constant-translation record; coefficients with
$q>d_i$ remain zero.  A shift $t\mapsto t+u$, $|u|\leq H\leq N$, transforms that
coefficient into $\sum_{v\geq q}\binom vq c_v u^{v-q}$, whose modulus is at most
$2^{d_i}M N^{d_i-q}$.  Subtracting the pivot multiplies this bound by at most two.
Adding a constant translation to a function record adds two such bounds at $q=0$.
Induction for the at most $R_0$ steps gives a fixed numerical multiple of the initial budget.
This proves every support-envelope and coefficient upper bound used by the analytic execution.

\paragraph{Lean state.}
Use an outer polynomial in $t$ with coefficients in a multivariate polynomial ring in the
shift variables.  A factor carries its origin index, conjugation flag, and constant-translation
polynomial.  Normalization groups \emph{function lists}; it does not discard duplicate factors.
The head list has length one.  Polynomial equality can be handled classically; computability
of comparisons between arbitrary real coefficients is not asserted.  The termination proof
uses the lexicographic well-founded relation, not an incorrect induction on maximum degree alone.

\subsubsection{The analytic update, with no assumed cube}
For fixed old shifts let
\[
 A=V^{-1}\int g_0(x)\mathbb E_{t\in[0,N]}
                   \prod_{i=1}^L g_i(x+Q_i(t))\,dx,
\]
where every block is $1$-bounded and has $L^2$ norm squared at most $V$.
Apply \cref{patch:signed-vdc} with $g=g_0$ and the remaining product as $F$.
Then translate $x$ by $-Q_p(t)$ in the correlation on its right side.
The two children of a factor $(g_i,Q_i)$ are exactly
\[
 (\overline{g_i},Q_i(t+u)-Q_p(t)),\qquad(g_i,Q_i(t)-Q_p(t)).
\]
Normalize constants as in \cref{patch:pet-update}.  The old spatial block was removed by the
Cauchy--Schwarz inequality; constant children form the new spatial block.  This proves
\begin{equation}\label{patch:pet-analytic-step}
 |A(\mathbf h)|^2\leq2\operatorname{Re}\mathbb E_{u\sim\kappa_H}
                       A'(\mathbf h,u)+2H/N.
\end{equation}
A product block is dominated by any one of its factors, so all required $L^2$ and $L^1$
bounds remain at most $V$.  The shifted head is a translated/conjugated copy of $f$.
Every joint integral is measurable; no numerical shift is selected during this symbolic recursion.

\subsubsection{The affine endpoint: removal order and cube invariant}
At a formal linear state, \eqref{patch:lead-invariant} says that the head slope and its differences
with every other slope are multiples of $e_m$.  Hence all the slopes are multiples of $e_m$.
Include the spatial block as a nuisance of slope zero, list all nuisances in their inherited
order $0,\ldots,s-1$, and place the protected block last, at index $s$.
For fixed old shifts the quantity has the exact form
\[
 A_0=V^{-1}\int\mathbb E_{t\in[0,N]}\prod_{i=0}^{s}g_i(x+a_it e_m)\,dx,
 \qquad g_s=\hbox{one translate/conjugate of }f.
\]
For $r\leq i\leq s$ define
\begin{align}
 G_{i,0}(x)&=g_i(x),\notag\\
 G_{i,r+1}(x;\mathbf u,u_r)&=G_{i,r}(x;\mathbf u)
     \overline{G_{i,r}(x+(a_i-a_r)u_re_m;\mathbf u)},\label{patch:affine-block}\\
 A_r(\mathbf u)&=V^{-1}\int\mathbb E_t
           \prod_{i=r}^sG_{i,r}(x+a_it e_m;\mathbf u)\,dx.\notag
\end{align}
Induction using a new Boolean bit proves the invariant
\begin{equation}\label{patch:affine-cube-invariant}
 G_{i,r}(x;\mathbf u)=\prod_{\omega\in\{0,1\}^r}\mathcal C^{|\omega|}
        g_i\left(x+\sum_{\nu<r}\omega_\nu(a_i-a_\nu)u_\nu e_m\right).
\end{equation}
At step $r$ translate by $-a_rt e_m$ and remove \emph{exactly} $G_{r,r}$ by
\cref{patch:signed-vdc}.  It is independent of $t$ after that translation.
Equation \eqref{patch:affine-block} identifies the resulting correlation with $A_{r+1}$, so
\[
 |A_r|^2\leq2\operatorname{Re}\mathbb E_{u_r\sim\kappa_H}A_{r+1}+2H/N.
\]
The zero vertex bounds $|G_{i,r}|$ by $|g_i|$, and $|A_r|\leq1$.
At $r=s$ translate away $a_st e_m$; the parameter average disappears and the result is
\begin{equation}\label{patch:affine-terminal}
 A_s(\mathbf u)=V^{-1}\int
     \prod_{\omega\in\{0,1\}^s}\mathcal C^{|\omega|}
         g_s\left(x+\sum_{\nu<s}\omega_\nu(a_s-a_\nu)u_\nu e_m\right)dx.
\end{equation}
Translation/conjugation invariance identifies its integrated power with that of $f$.
The raw cube in \eqref{patch:affine-terminal} makes sense even at a parameter where a slope
difference evaluates to zero.  We only require nonzero gaps later, for changing variables.
Thus no exceptional set is hidden in the analytic recursion.

\begin{lemma}[Explicit accumulation of the removal losses]\label{patch:cs-loss}
Suppose a chain of $T$ such steps has $|A_r|\leq1$, probability shift measures, and error
$2H/N$ at each step.  If its final signed average is a nonnegative cube power $Q\leq1$, then
\[
 |A_0|^{2^T}\leq c_T\bigl(Q+T H/N\bigr),\qquad c_T=2^{2^T-1}.
\]
The same assertion holds with all earlier shift parameters retained and averaged.
\end{lemma}
\begin{proof}
Set $M_r=\mathbb E|A_r|$ except at the last stage, where use $Q=\operatorname{Re}\mathbb E A_T$.
Jensen gives $M_r^2\leq2M_{r+1}+2H/N$, with $Q$ in the last line.
For $\phi(x)=x^2/2$, $x,y\in[0,1]$ imply
$|\phi(x)-\phi(y)|\leq|x-y|$.  Inductively
$\phi^{\circ r}(M_0)\leq M_r+rH/N$, and at the last stage use $Q$.
For a completely scalar induction, split into $\phi^{\circ r}(M_0)\leq M_r$ or its reverse
and use $(x^2-y^2)/2=(x-y)(x+y)/2$ on $[0,1]$.
Since $\phi^{\circ T}(x)=x^{2^T}/c_T$, the result follows.
The last signed average is essential; it is not replaced by absolute cube correlations.
\end{proof}

\subsubsection{Real sublevel exclusion and the common-scale cube}
At the terminal stage every head gap, including the gap with the spatial slope zero, has form
\[
 a\,p_i(\mathbf h),\qquad p_i\in\mathbb Z[\mathbf h]\setminus\{0\},
\]
where $p_i$ is multilinear homogeneous of degree $D-1\leq2$, with a coefficient bound fixed
by the recursion above.  All old parameter shifts have the same radius $H$.

\begin{lemma}[The sublevel estimate actually needed here]\label{patch:multiaffine-sublevel}
Let $p\ne0$ be a multilinear integer polynomial of degree at most two.  If the coordinates
of $u$ are independently distributed with density $\kappa_1$, then, for $0<\varepsilon<1$,
\[
 \Pr\{|p(u)|\leq\varepsilon\}\leq4\sqrt\varepsilon.
\]
\end{lemma}
\begin{proof}
A nonzero constant has modulus at least one.  For a nonconstant linear polynomial, fix all
but a coordinate with a nonzero integer coefficient.  The sublevel interval has length at
most $2\varepsilon$, and $\kappa_1\leq1$.
If a quadratic term is present, write $p=A(u')u_r+B(u')$ where $A$ is linear and has a
nonzero integer coefficient.  The probability $|A|\leq\sqrt\varepsilon$ is at most
$2\sqrt\varepsilon$ by the linear case.  On its complement, the conditional sublevel
probability in $u_r$ is at most $2\varepsilon/|A|\leq2\sqrt\varepsilon$.
Integrate the remaining probability variables.  No dimension-dependent inverse theorem is used.
\end{proof}

Let $T_0$ bound the total number of parameter and affine removals; for example use
$T_0=4+R_0+L_1$.  If the original normalized lower bound is $\lambda\leq1$, choose
\[
 H/N=\min\{1/32,\lambda^{2^{T_0}}/(4T_0c_{T_0})\}>0.
\]
Take $\lambda$ to be the fixed lower bound returned from the input budget, rather than
the measured value of the particular correlation.  This choice of $H/N$ has both a power
lower bound and a power upper bound by the budget rules.  The loss lemma gives an averaged raw
terminal cube at least $\kappa=\lambda^{2^{T_0}}/(2c_{T_0})$ after harmless weakening.
For $s$ terminal gaps, take $\varepsilon=(\kappa/(16s))^2$.
Homogeneity gives $p_i(Hu)=H^{D-1}p_i(u)$.
By \cref{patch:multiaffine-sublevel}, the union of all bad gap sets has probability at most
$\kappa/4$.  Since the raw cube is between zero and one, there are old shifts outside this
union with raw cube at least $\kappa/2$.  At those shifts
\[
 |a|H^{D-1}\varepsilon\leq |a p_i(\mathbf h)|
       \leq |a|H^{D-1}M_1,
\]
where $M_1$ is the integer coefficient $\ell^1$ bound.
Changing the new affine increments by $z_i=a p_i(\mathbf h)u_i$ therefore produces
Fej\'er radii $L_i=|a p_i(\mathbf h)|H$ with
\begin{equation}\label{patch:pet-physical-radii}
 \delta^{O(1)}N^D\leq L_i\leq\delta^{-O(1)}N^D.
\end{equation}
Here $D$ is the protected polynomial degree, not the sum of ambient degrees.

\begin{lemma}[Mixed radii to one radius, with positivity before comparison]\label{patch:uniformize}
Suppose $f$ is $1$-bounded on a box of volume $V$ and $j$-th side length $S$.
If $L_i>0$ and $L\geq2\max_iL_i$, then
\[
 (Q_{\mathbf L,V}^{j}(f))^2
 \leq\left(1+\frac LS\right)\prod_i\frac{2L}{L_i}\ Q_{s+1,L,V}^{j}(f).
\]
\end{lemma}
\begin{proof}
For fixed $z_1,\ldots,z_s$ put $u_z=C_s(f;\cdot,z)$.
The zero vertex supports $u_z$ on the original box, so
\[
 |V^{-1}\!\int u_z|^2\leq(1+L/S)\,P(z),\qquad
 P(z)=V^{-1}\|T_{L,j}u_z\|_2^2\geq0.
\]
Since $L_i\leq L/2$, pointwise
$\kappa_{L_i}(z_i)\leq(2L/L_i)\kappa_L(z_i)$.
Apply this comparison only to the nonnegative $P$, after Jensen in the first display.
Square expansion of $P$ appends one cube coordinate by \eqref{patch:cube-succ}; all $s+1$
increments now have the same Fej\'er density.  This proves the identity and inequality claimed.
\end{proof}
Choose a deterministic $L=u_C(\delta)N^D$ exceeding twice the upper bounds in
\eqref{patch:pet-physical-radii} and comparable, up to a power of $\delta$, to the protected
support side length.  Every factor in \cref{patch:uniformize} is power-controlled.
If needed, pad the resulting order to a single fixed order using \eqref{patch:raise-order}.
This proves highest-active-degree uniformity control with polynomial losses, including the
shifted families \eqref{patch:initial-shift-family}; its protected factor is still the original $f$.

\subsubsection{Polynomial phases: the four initial steps are also explicit}
Suppose the correlation also contains $\ee{p_x(t)}$, with $p_x$ a measurable-in-$x$ polynomial
in $t$ of degree at most three.  Its coefficients need not be bounded.
Before any polynomial-dependent translation of $x$, apply \cref{patch:signed-vdc} four times,
retaining $x$ as a fixed outer variable.  After $r$ steps the parameter-dependent part is
\[
 \prod_{i=1}^{m}\prod_{\omega\in\{0,1\}^r}\mathcal C^{|\omega|}
     f_i(x-P_i(t+\omega\cdot\mathbf h)e_i)\;
 \ee{\sum_{\omega\in\{0,1\}^r}(-1)^{|\omega|}p_x(t+\omega\cdot\mathbf h)}.
\]
The initial step removes the original spatial factor.  At later steps use the indicator of
a fixed enlarged spatial box: the zero-shift vertex supports the integrand there for
$t\in[0,N]$.  Polynomial coefficient bounds and $4H\leq N/8$ give a single power-controlled
volume for this box.  The indicator is removed by the same Cauchy--Schwarz step.
At $r=4$ the phase is identically one, by the fourth finite-difference identity for cubics.
Normalize the constant terms of the remaining polynomial shifts and run the preceding PET
producer, retaining all four old shift variables in its averages and sublevel exclusions.

There is one degenerate initial case to treat separately.  If the highest active degree is
one, there is only one active input.  After the four steps, translating $x$ by its affine
$P_1(t)$ removes $t$ and gives the raw order-four cube of $f_1$ with increments
$\operatorname{lc}(P_1)h_i e_1$ directly.  Keep the signed average at this last step.
Its radius is $|\operatorname{lc}(P_1)|H>0$ and is comparable to $N$ up to powers of $\delta$.
Uniformize and pad as above.  Do \emph{not} group these equal-slope head copies and call the
result a singleton head.

\begin{proposition}[Highest-active-input control, also with measurable polynomial phases]
\label{patch:highest-control}
Fix an increasing degree list in $\{1,2,3\}$, a support/coefficient budget, and a number of
active inputs $m\geq1$.  A power lower bound for
\[
 \left|\int f_0(x)\mathbb E_{t\in[0,N]}
           \prod_{i=1}^{m}f_i(x-P_i(t)e_i)\,\ee{p_x(t)}\,dx\right|
\]
with $1$-bounded inputs, admissible polynomials, budgeted support, and $\deg_t p_x\leq3$
implies a power lower bound for a fixed-order directional uniformity norm of $f_m$,
at a radius comparable to $N^{d_m}$ up to fixed powers of $\delta$.
The order and all budgets are independent of the measurable phase coefficients and of the
input functions.  Passive coordinates, if present, remain outer integration variables.
\end{proposition}
\begin{proof}
The four phase steps, the normal-state update with invariant \eqref{patch:lead-invariant},
the explicit pivot recursion, \eqref{patch:affine-terminal}, the loss lemma, the sublevel
exclusion, and \cref{patch:uniformize} give the result.  Each ingredient has been proved above.
The degree-one branch is the direct cube calculation just given.  Dividing by the chosen
spatial volume and finally converting to the stated box normalization costs only a budgeted
factor; the exact ambient factor remains $N^{\sum d_i}$.
\end{proof}
\paragraph{Lean nodes.}
\mbox{}\par
{\small\raggedright
\pid{NormalPETState}; \pid{ProtectedLeadingInvariant}; \pid{petUpdate_protected};
\pid{petPivot_type_decreases}; \pid{petIteration_bound}; \pid{petStep_signed};
\pid{affineBlock_cube}; \pid{csChain_power_bound}; \pid{multiaffine_sublevel};
\pid{mixedCube_uniformize}; \pid{phaseStrip_four}; \pid{highestActive_uniformity}.
\par}
The proposition is not an extra hypothesis in any later theorem.


\subsection{Fourier selection and the two difference lemmas}\label{patch:auxiliary-differences}
These are the real versions of the preparation tools in \cite[\S4.5.1]{KoszEtAl}.
Their Cauchy--Schwarz steps are expanded here using \cref{patch:cylinder}.
For a scalar function supported on an interval $J$ of length $S>0$, define the paired cube
\[
 D_{a,b}^{r}f(x)=\prod_{\omega\in\{0,1\}^r}\mathcal C^{|\omega|}
       f\left(x+\sum_{i=1}^r[\omega_i a_i+(1-\omega_i)b_i]\right),
 \quad a,b\in[0,H]^r.
\]
After translation in $x$, averaging this expression over $a,b$ is the same as averaging
$C_r$ against $r$ independent Fej\'er densities.  All scalar normalizers in this subsection
are written as $S$; changing to a larger interval costs its explicit length ratio.

\begin{lemma}[One-dimensional $U^2$ inverse estimate and measurable frequencies]
\label{patch:u2-fourier-selection}
If $|f|\leq1$ and $f$ is supported in $J$, then
\[
 Q_{2,H,S}(f)\leq H^{-2}\|\widehat f\|_\infty^2.
\]
For a jointly measurable bounded family with uniformly bounded support, a power lower bound
for these powers on a measurable set of parameters gives a measurable choice of frequencies
with Fourier coefficients of power size.  A finite frequency set is not assumed.
\end{lemma}
\begin{proof}
Put $g_a(x)=f(x)\overline{f(x+a)}$.  By \cref{patch:fejer-squares},
$SQ_{2,H,S}=\int\kappa_H(a)\|T_Hg_a\|_2^2da$.
Use $\kappa_H\leq H^{-1}$ against this nonnegative integrand and apply Plancherel in $x$.
For each $\xi$, a second Plancherel identity, this time in $a$, gives
\[
 \int|\widehat{g_a}(\xi)|^2da
    =\int|\widehat f(\zeta+\xi)|^2|\widehat f(\zeta)|^2d\zeta
    \leq\|\widehat f\|_\infty^2\|f\|_2^2.
\]
The multiplier of $T_H$ is the Fourier transform of $H^{-1}\mathbf1_{[0,H]}$, whose squared
$L^2$ norm is $H^{-1}$.  Hence $SQ_2\leq H^{-2}\|\widehat f\|_\infty^2\|f\|_2^2$;
use $\|f\|_2^2\leq S$.
For selection enumerate $\mathbb Q$.  The Fourier transform of each integrable section is
continuous in frequency, and its value at each rational is measurable in the parameter by
Fubini.  Where the supremum exceeds a threshold $a$, choose the first rational at which the
modulus exceeds $a/2$.  The disjoint least-index sets are measurable and cover the required
parameter set.  Outside it use zero.  Popularity supplies a uniform positive threshold
before selection.  This proves the claimed jointly measurable choice and does not infer a
uniform finite set from continuity alone.
\end{proof}

\begin{lemma}[Dual--difference interchange, explicit cylinder proof]\label{patch:dual-difference}
Let $F(x)=\mathbb E_tF_t(x)$, where the $F_t$ are jointly measurable, $1$-bounded and supported
in a common interval $J$.  Let $E\subseteq[0,H]^{2r}$ and $\varphi(a,b)$ be measurable, $r\geq1$.
Put
\[
 \theta=\mathbb E_{a,b\in[0,H]^r}\mathbf1_E(a,b)
           \left|S^{-1}\int D_{a,b}^{r}F(x)\ee{\varphi(a,b)x}\,dx\right|^2.
\]
Write $J_r=J-[0,rH]$, $T=|J_r|/S$, and
\[
 \Box^r E=\{(a,b,c): (a,b^\omega)\in E\text{ for every }\omega\in\{0,1\}^r\},
 \quad b_i^\omega=(1-\omega_i)b_i+\omega_i c_i.
\]
Then, with $G_{b,c}(x)=\mathbb E_tD_{c,b}^{r}F_t(x)$ and
$\psi(a,b,c)=\sum_\omega(-1)^{|\omega|}\varphi(a,b^\omega)$,
\begin{equation}\label{patch:dual-diff-bound}
 \mathbb E_{a,b,c}\mathbf1_{\Box^r E}
       \left|S^{-1}\int G_{b,c}(x)\ee{\psi(a,b,c)x}\,dx\right|^2
 \geq\theta^{2^r}T^{-2(2^r-1)}.
\end{equation}
\end{lemma}
\begin{proof}
Expand the square defining $\theta$ using two spatial variables $x,z$ and expand every copy
of $F$ using independent parameters $t_{\omega,0},t_{\omega,1}$.
The vertex $\omega=0$ contributes
\[
 v=\mathbf1_E(a,b)F_{t_{0,0}}(x+\textstyle\sum b_i)
       \overline{F_{t_{0,1}}(z+\textstyle\sum b_i)}\,
       \ee{\varphi(a,b)(x-z)}.
\]
Partition the other Boolean vertices by their first nonzero coordinate:
$\Omega_i=\{\omega:\omega_1=\cdots=\omega_{i-1}=0,\ \omega_i=1\}$.
The product of their two spatial factors for $\omega\in\Omega_i$ is a $1$-bounded nuisance
$u_i$ independent of $b_i$.  Normalize $x,z$ on $J_r$, and normalize all parameter variables
by their probability measures.  The initial average is $\theta/T^2$.
Apply \cref{patch:cylinder}, doubling $b_1,\ldots,b_r$ in that order.
At stage $i$ it removes exactly the block indexed by $\Omega_i$ and hence all its unused
$t$ parameters.  The only parameters surviving at the endpoint are $t_{0,0},t_{0,1}$.
The indicator becomes the product over all corners, namely $\mathbf1_{\Box^r E}$;
the phase becomes $\ee{\psi(x-z)}$; the two function products become
$D_{c,b}^{r}F_{t_{0,0}}(x)$ and its conjugate at $z$.
Integrating $t_{0,0},t_{0,1},x,z$ is exactly the squared coefficient of $G_{b,c}$.
The support remains in $J_r$, so these spatial integrals may be written over $\mathbb R$.
Multiply the cylinder inequality by $T^2$ to obtain \eqref{patch:dual-diff-bound}.
\end{proof}

\begin{lemma}[Removal of missing-coordinate phases]\label{patch:missing-phase}
Let $f$ be as above and let $\varphi_i(a,b)$ be independent of $b_i$, $1\leq i\leq r$.
If
\[
 \theta=\mathbb E_{a,b}
       \left|S^{-1}\int D_{a,b}^{r}f(x)\ee{(\sum_i\varphi_i(a,b))x}\,dx\right|^2,
\]
then
\begin{equation}\label{patch:missing-phase-bound}
 Q_{r+1,H,S}(f)\geq
       \frac{\theta^{2^r}}{T^{2(2^r-1)}(1+H/S)},\qquad T=(S+rH)/S.
\end{equation}
\end{lemma}
\begin{proof}
Expand the square using $x,z\in J_r$.
Use as head $v=f(x+\sum b_i)\overline{f(z+\sum b_i)}$.
For nuisance $u_i$, multiply the vertex factors indexed by $\Omega_i$ from the preceding
proof by $\ee{\varphi_i(a,b)(x-z)}$.  This block is independent of $b_i$.
The same cylinder induction gives
\[
 R:=\mathbb E_{b,c}\left|S^{-1}\int D_{c,b}^{r}f(x)\,dx\right|^2
       \geq\theta^{2^r}T^{-2(2^r-1)}.
\]
For each $b,c$, $u=D_{c,b}^{r}f$ is supported in a translate of $J$, by its zero vertex.
Its full integral equals $\int T_Hu$, and the latter function has support length at most $S+H$.
Thus $|S^{-1}\int u|^2\leq(1+H/S)S^{-1}\|T_Hu\|_2^2$.
Average in $b,c$ and expand this last square.  It is exactly the paired cube with $r+1$
uniform pairs in $[0,H]$, hence $Q_{r+1,H,S}$.  This proves \eqref{patch:missing-phase-bound}.
\end{proof}
These two lemmas are useful even when the support/shift ratios are large: the displayed
ratios are then retained and bounded by the budget rules.  No comparison of complex
integrands by positive kernels is implicit.

\begin{lemma}[A measurable dummy phase with a small exceptional set]\label{patch:dummy-phase}
For $r\geq1$, $H,M>0$, and $0<\epsilon<1$, there is a finite-valued measurable function
$\Phi(a,b)$ on $[0,H]^{2r}$ such that its alternating sum
$\Psi(a,b,c)=\sum_\omega(-1)^{|\omega|}\Phi(a,b^\omega)$ satisfies $|\Psi|>M$ except on
a set of relative measure at most $\epsilon$ in $[0,H]^{3r}$.
\end{lemma}
\begin{proof}
Choose $n$ with $r2^{-n}\leq\epsilon$ and partition each $[0,H]$ into $2^n$ half-open
intervals, with the final endpoint assigned to the last cell.
Enumerate the resulting $r$-dimensional cells by distinct nonnegative integers $\iota$.
On a cell of the $b$ variable put $\Phi(a,b)=2M\,2^{\iota(b)}$.
Except when some $b_i,c_i$ are in the same one-dimensional cell, the $2^r$ corners $b^\omega$
are in distinct cells.  A signed sum of distinct powers of two has absolute value at least
one: its largest term exceeds the sum of all smaller possible terms.
Thus $|\Psi|\geq2M$ there.  The union of the $r$ bad coordinate events has relative measure
at most $r2^{-n}$.  Boundaries are null.  This grid resolves the purpose-built function
$\Phi$, not the level sets of an arbitrary input.  The possibly enormous phase values are
permitted: later phase mappings are arbitrary measurable functions, and their exponentials
still have modulus one.
\end{proof}

\subsection{Conditional degree lowering with all hypotheses exposed}\label{patch:conditional-dl}
Let $\mathbb R^0$ be a point with measure one.  For fixed $k,l,m$ with $1\leq m\leq l\leq k$,
define $\mathsf{MA}(m,l)$ to be the following budgeted assertion.  For every input budget
there is an output budget such that, for admissible $P_i$, $1$-bounded $g_0,\ldots,g_{m-1}$
on $\mathbb R^{m-1}$, measurable $\zeta_m,\ldots,\zeta_l$, and fixed
$\theta_{l+1},\ldots,\theta_k$, a power lower bound for
\begin{equation}\label{patch:ma-correlation}
 \left|\int g_0(y)\mathbb E_t\prod_{i<m}g_i(y-P_i(t)e_i)
       \ee{\sum_{i=m}^l\zeta_i(y)P_i(t)+\sum_{i>l}\theta_iP_i(t)}\,dy\right|
\end{equation}
in units of $N^{D_{m-1}}$ implies that the canonical set
\begin{equation}\label{patch:ma-set}
 Y=\left\{y\text{ in the support envelope}:\
      \sum_{i=m}^lN^{d_i}|\zeta_i(y)|+\sum_{i>l}N^{d_i}|\theta_i|
                           \leq u_C(\delta)\right\}
\end{equation}
has measure at least $\ell_C(\delta)N^{D_{m-1}}$.
There is no bound on the phase mappings before this conclusion.  In the real case no
integer denominator occurs.  The output budget is independent of the phase values.

\begin{lemma}[Conditional structured degree lowering]\label{patch:conditional-degree}
Assume the theorem $\mathsf{MA}(m,l)$ for every budget.  For $s\geq3$ and admissible data set
\begin{align*}
 F(x)&=\mathbb E_t f_0(x+P_m(t)e_m)
           \prod_{i<m}f_i(x-P_i(t)e_i+P_m(t)e_m)\\
 &\hspace{13mm}\cdot
 \ee{\sum_{i=m+1}^l\xi_i(x+P_m(t)e_m)P_i(t)+\sum_{i>l}\eta_iP_i(t)}.
\end{align*}
Here the $f_i$ are $1$-bounded with budgeted support, the $\xi_i$ are measurable, and the
$\eta_i$ are fixed.  If $H$ is comparable to $N^{d_m}$ up to budgeted powers and the
order-$s$ uniformity norm of $F$ in direction $e_m$ is power-large, then its order-$(s-1)$
norm at the \emph{same} $H$ is power-large.  Enlarge the support normalizer to contain $F$;
its volume ratio is explicitly power-controlled.
\end{lemma}
\begin{proof}
Write $r=s-2$, $x=(y,z)$ with $y=x_{<m}$ and $z=x_m$.
The following steps specify every selection and every use of the assumed major-arc theorem.

\paragraph{1. Fourier coefficients on popular sections.}
Expand the first $r$ paired differences in the order-$s$ power.  Fubini identifies the
remainder with a one-dimensional $U^2$ power on each $y$ section.
Use \cref{patch:budgets,patch:u2-fourier-selection}.  There is a measurable power-large set
of $y$ and, on each such section, a measurable power-large set $E_y\subseteq[0,H]^{2r}$
on which
\[
 \left|S^{-1}\int D_{a,b}^{r}F_y(z)\ee{\varphi_y(a,b)z}\,dz\right|\geq\delta^{O(1)}.
\]
Here $S$ is the length of one fixed enlarged section interval $J$, comparable to $N^{d_m}$
up to powers.  The frequency $\varphi_y(a,b)$ is the least-rational selection described above,
so it is jointly measurable.  Translated supports of the differences are all contained in
$J-[0,rH]$ after one further budgeted enlargement if necessary.

\paragraph{2. Interchange the average and the differences.}
Write $F_y=\mathbb E_tF_{t,y}$ and apply \cref{patch:dual-difference} on each good $y$.
Define
\[
 G_{y,b,c}(z)=\mathbb E_tD_{c,b}^{r}F_{t,y}(z),\qquad
 \psi_y(a,b,c)=\sum_\omega(-1)^{|\omega|}\varphi_y(a,b^\omega).
\]
Popularity gives a canonical measurable set $Z$ of pairs $(y,k)$, $k=(a,b,c)$, on which
$(a,b^\omega)\in E_y$ for every corner and
$|\int G_{y,b,c}(z)\ee{\psi_y(k)z}dz|\geq\delta^{O(1)}N^{d_m}$.
The measure of $Z$ is power-large in the units $N^{D_{m-1}}H^{3r}$.

\paragraph{3. Interchange the section and cube parameters.}
Fubini and popularity produce a measurable set $K\subseteq[0,H]^{3r}$ of power-large relative
measure such that each $k\in K$ has a power-large section $Z_k$ in the $y$ variables.
Every later set can be defined by explicit thresholds on these joint measurable functions;
no arbitrary sectionwise choice of a set is made.

\paragraph{4. Extend by a dummy phase.}
On $Z$ retain $\psi_y(k)$.  Off $Z$ replace it by the phase $\Psi(k)$ from
\cref{patch:dummy-phase}.  Denote the result by $\psi_y^*(k)$.
Choose its threshold $M$ larger than the frequency bound returned in Step 5, and its exceptional
relative measure smaller than one quarter of the power mass used in Step 6.
These choices are legitimate: Step 5's constants do not depend on phase values.
The modified phase is not asserted to be a cube derivative off $Z$.
For every $k\in K$,
\[
 \int\left|\int G_{y,b,c}(z)\ee{\psi_y^*(k)z}dz\right|dy
                    \geq\delta^{O(1)}N^{D_m}.
\]

\paragraph{5. Apply precisely $\mathsf{MA}(m,l)$.}
Choose a measurable unit phase in $y$ to turn the inner coefficient into its modulus.
Expand $G$, distribute the paired derivative over the factors, and set the new spatial
variable equal to $x+P_m(t)e_m$.  The fixed $\eta$ phase cancels because $r\geq1$.
The resulting integral is
\[
 \int f'_0(x)\mathbb E_t\prod_{i<m}D_{c,b}^{r}f_i(x-P_i(t)e_i)
           \ee{\sum_{i=m}^l\zeta'_i(x)P_i(t)}\,dx,
\]
where
\[
 \zeta'_m(y,z)=-\psi_y^*(k),\qquad
 \zeta'_i(x)=\sum_\omega(-1)^{|\omega|}\xi_i(x+\textstyle\sum_j b_j^\omega e_m)
       \quad(i>m),
\]
and $f'_0$ is the product of the chosen $y$ phase, $D_{c,b}^{r}f_0(x)$, and
$\ee{\psi_y^*(k)x_m}$.  It is $1$-bounded and supported in a budgeted box.
Fix a $z$ section carrying a power-large correlation.  This is exactly
\eqref{patch:ma-correlation}, in $m-1$ variables, with the same $(m,l)$ and zero fixed higher
phase.  The assumed $\mathsf{MA}(m,l)$ yields a power-large set of $y$ on which
$|\psi_y^*(k)|\leq u_{C_6}(\delta)N^{-d_m}$.
Take as this set the full threshold set.  It is jointly measurable in $(y,k)$, whether or
not the intermediate selected $z$ depends measurably on $k$.

\paragraph{6. Recover the original cube phase.}
Choose the dummy threshold in Step 4 greater than $u_{C_6}(\delta)N^{-d_m}$.
Outside its small exceptional set in $k$, every low-frequency point just obtained lies in
$Z$, where $\psi_y^*=\psi_y$.
Removing the exceptional set and applying Fubini/popularity again leaves power-many $y$,
each with power-many $k=(a,b,c)$ for which the original alternating sum is small and all
corners belong to $E_y$.  All losses are bounded by the explicit discard rule in
\cref{patch:budgets}.

\paragraph{7. Quantize and partition the nonzero vertices.}
Partition the interval of possible $\psi_y(k)$ values into cells of width $\tau N^{-d_m}$.
Choose a common cell carrying power mass, with center $\zeta$; the number of cells is
power-controlled.  Choose $\tau>0$ small enough that
$2\pi(\sup_{z\in J-[0,rH]}|z|)\tau N^{-d_m}$ is at most a quarter of the normalized Fourier
coefficient lower bound in Step 1.  This choice is made before the loss from selecting a cell.
For $\Omega_i=\{\omega:\omega_{<i}=0,\omega_i=1\}$ define
\[
 \alpha_1=\zeta-\sum_{\omega\in\Omega_1}(-1)^{|\omega|}\varphi_y(a,b^\omega),
 \qquad
 \alpha_i=-\sum_{\omega\in\Omega_i}(-1)^{|\omega|}\varphi_y(a,b^\omega)\ (i>1).
\]
Each $\alpha_i$ is independent of $b_i$, and on the retained set
\[
 \left|\varphi_y(a,b)-\sum_i\alpha_i(y,a,b,c)\right|
       =|\psi_y(a,b,c)-\zeta|\leq\tau N^{-d_m}.
\]
This is an exact partition of all nonzero Boolean vertices, not an asserted separation of
variables.

\paragraph{8. Remove the phases and integrate.}
For each good $y$, fix one $c$ for which the retained $(a,b)$ section has power-large measure.
On it, replacing the Fourier phase in Step 1 by $\sum_i\alpha_i$ loses at most a quarter of
its lower bound, by the choice of $\tau$ and $|\ee u-\ee v|\leq2\pi|u-v|$.
Enlarge the average to all $(a,b)$ only after taking the squared modulus, which is nonnegative.
Apply \cref{patch:missing-phase}.  Since $r+1=s-1$, it gives a power lower bound for the
order-$(s-1)$ norm of $F_y$ at the same $H$.
There is no need to choose $c$ measurably: the resulting norm inequality is an intrinsic
measurable statement about $F_y$ and holds for every good $y$.
Integrate over the already measurable good set of $y$.  Section normalization ratios,
shift/support ratios, and the fixed roots involved are all handled by \cref{patch:budgets}.
\end{proof}
\paragraph{Dependency warning built into the implementation.}
\pid{conditionalDegreeLowering} is a theorem with an internal premise
\pid{MajorArcProperty m l}; it is proved \emph{before} the major-arc induction.
The next subsection supplies that premise by induction.  The final every-input theorem
will have no such premise.


\subsection{The major-arc induction is noncircular}\label{patch:major-induction}
\begin{theorem}[Real major-arc property]\label{patch:major-arc}
For every increasing degree list contained in $\{1,2,3\}$ and every $1\leq m\leq l\leq k$,
the property $\mathsf{MA}(m,l)$ defined in \eqref{patch:ma-correlation}--\eqref{patch:ma-set}
holds.  All output budgets depend only on the input budget and the degree list.
\end{theorem}
\begin{proof}
Fix $k,l$ and induct on $m=1,\ldots,l$, with the induction hypothesis quantified over
\emph{every} support/coefficient/correlation budget.

\paragraph{Base $m=1$.}
The space $\mathbb R^0$ has one point and measure one.  Since $|g_0(0)|\leq1$, a lower
bound for \eqref{patch:ma-correlation} gives the same lower bound for the scalar exponential
average.  Apply \cref{patch:triangular}.  The canonical set in \eqref{patch:ma-set} is the
whole point space after increasing its budget, and its measure is one.

\paragraph{Induction step: replace the last active input by its exact adjoint.}
Suppose $\mathsf{MA}(m,l)$ is proved, and consider the correlation for $\mathsf{MA}(m+1,l)$:
\[
 \Gamma=\int g_0(x)\mathbb E_t\prod_{i\leq m}g_i(x-P_i(t)e_i)\,
          \ee{p_x(t)}\,dx,\qquad
 p_x(t)=\sum_{i=m+1}^l\zeta_i(x)P_i(t)+\sum_{i>l}\theta_iP_i(t).
\]
Here $x\in\mathbb R^m$.  Define, with all conjugations shown,
\begin{equation}\label{patch:ma-adjoint}
 F(y)=\mathbb E_t\overline{g_0(y+P_m(t)e_m)}
     \prod_{i<m}\overline{g_i(y+P_m(t)e_m-P_i(t)e_i)}\,
               \ee{-p_{y+P_m(t)e_m}(t)}.
\end{equation}
Translation in $x$ gives $\Gamma=\langle g_m,F\rangle$.
If the $g_i$ have squared $L^2$ norms bounded by $V$, then
\[
 \|F\|_2^2\geq |\Gamma|^2/V.
\]
The function $F$ is $1$-bounded and is supported in a fixed enlarged box, because its first
factor is a translate of $g_0$.  Replacing $g_m$ by $F$ in $\Gamma$ produces exactly
$\langle F,F\rangle=\|F\|_2^2$.  Thus the replaced correlation remains power-large in
units $N^{D_m}$.

\paragraph{Obtain uniformity, using only the smaller major-arc property.}
The phase in the replaced correlation is still $p_x(t)$, a polynomial in $t$ of degree at
most three with measurable coefficients.  Apply \cref{patch:highest-control}, whose proof
is independent of every major-arc assertion.  It gives power-large order-$S_0$ uniformity
of $F$ in direction $e_m$, at a radius $H$ comparable to $N^{d_m}$ up to powers.
The order $S_0$ is bounded independently of all functions and parameters; pad it to a fixed
order if necessary using \eqref{patch:raise-order}.
The function \eqref{patch:ma-adjoint} has exactly the form in
\cref{patch:conditional-degree}, with conjugated $g_i$, mappings $\xi_i=-\zeta_i$, and
constants $\eta_i=-\theta_i$.
Use the induction hypothesis $\mathsf{MA}(m,l)$ in that conditional lemma, and lower the
order $S_0-2$ times, at the same $H$.  If $S_0=2$ there is no lowering step.
This obtains a power-large directional $U^2$ power of $F$.

\paragraph{Fourier-select and expose the smaller pattern.}
By Fubini, popularity, and \cref{patch:u2-fourier-selection}, there is a measurable mapping
$\lambda:\mathbb R^{m-1}\to\mathbb R$ such that
\[
 \int\left|\int F(y,z)\ee{\lambda(y)z}\,dz\right|dy
                      \geq\delta^{O(1)}N^{D_m}.
\]
Take a measurable scalar $c(y)$ of modulus at most one that turns the inner coefficient
into its modulus.  Expand $F$ and put the new variable $x$ equal to the old variable plus
$P_m(t)e_m$.  The resulting power-large scalar is exactly
\begin{equation}\label{patch:ma-smaller-pattern}
 \int \overline{g_0(x)}c(x_{<m})\ee{\lambda(x_{<m})x_m}
 \mathbb E_t\prod_{i<m}\overline{g_i(x-P_i(t)e_i)}
 \ee{-\lambda(x_{<m})P_m(t)-p_x(t)}\,dx.
\end{equation}
The $x_m$ integration is restricted by $g_0$ to an interval of power-controlled length
in units $N^{d_m}$.  Each section has absolute value at most a power-controlled multiple
of $N^{D_{m-1}}$.  Since the integral of the absolute values of the sections is at least
the modulus of \eqref{patch:ma-smaller-pattern}, popularity gives a measurable set $Z_m$
of $x_m$ values with measure at least $\delta^{O(1)}N^{d_m}$, on every one of which the
section correlation is at least $\delta^{O(1)}N^{D_{m-1}}$.
Use a fixed threshold for all these sections; do not feed their individual sizes into the
induction hypothesis.

For each $z\in Z_m$, the section of \eqref{patch:ma-smaller-pattern} is an
$\mathsf{MA}(m,l)$ correlation, with variable phases
\[
 -\lambda(y),\quad -\zeta_{m+1}(y,z),\ldots,-\zeta_l(y,z)
\]
and fixed higher phases $-\theta_i$.
Apply the induction hypothesis with one common budget.  Drop the nonnegative term
$N^{d_m}|\lambda(y)|$ from its conclusion.  It follows that the section of the canonical set
\[
 Y=\left\{(y,z)\text{ in the enlarged box}:\
       \sum_{i=m+1}^lN^{d_i}|\zeta_i(y,z)|+
       \sum_{i>l}N^{d_i}|\theta_i|\leq u_C(\delta)\right\}
\]
has measure at least $\delta^{O(1)}N^{D_{m-1}}$ for each $z\in Z_m$.
This set is jointly measurable, with no sectionwise set-selection issue.
Fubini yields $|Y|\geq\delta^{O(1)}N^{D_m}$, which is $\mathsf{MA}(m+1,l)$.

All losses in a step are products, fixed powers, or finite-order applications of the budget
rules.  Choose the input scale threshold large enough that every invoked theorem is
applicable.  There are at most $k-1\leq2$ induction steps, although each step invokes the
previously bounded number of degree-lowering iterations.  This proves the claimed uniform
structural budgets.
\end{proof}

The dependency at stage $m$ is precisely
\[
 \mathsf{MA}(m,l)
 \ \Longrightarrow\ \mathsf{ConditionalDL}(m,l)
 \ \Longrightarrow\ \mathsf{MA}(m+1,l),
\]
with highest-active uniformity proved before this induction.  In particular, neither the
conclusion at $m+1$ nor the final inverse theorem is used as an induction hypothesis.

\begin{corollary}[Unconditional degree lowering for structured adjoints]
\label{patch:structured-degree}
Let $2\leq l\leq k\leq3$ and fix admissible polynomials and bounded, budget-supported inputs.
Let $F_l$ be the $l$-th adjoint of the modulated average with active inputs $1,\ldots,l$
and a fixed phase $\ee{\sum_{i>l}\zeta_iP_i(t)}$.  If its directional order-$s$ norm,
$s\geq3$, is power-large at $H\asymp\delta^{O(1)}N^{d_l}$, then its order-$(s-1)$ norm
at the same $H$ is power-large, with uniform budgets.
\end{corollary}
\begin{proof}
Keep the passive coordinates $x_{>l}$ fixed.  The active $l$-variable function is exactly
the function in \cref{patch:conditional-degree} with $m=l$ and fixed higher phase,
up to the explicitly harmless conjugations.  Popularity first restricts to passive sections
with power-large order-$s$ power.  Apply that lemma using the now proved
$\mathsf{MA}(l,l)$, with the same budgets on all selected sections, and integrate the
resulting nonnegative order-$(s-1)$ powers.  The passive-section volume has its exact factor
$N^{D-D_l}$; support ratios are handled by \cref{patch:budgets}.
\end{proof}
This corollary is the structured statement corresponding to \cite[Lemma 4.70]{KoszEtAl}.
It makes no assertion that a generic function with large high-order uniformity has large
low-frequency energy or admits degree lowering.
\paragraph{Lean nodes.}
\mbox{}\par
{\small\raggedright
\pid{majorArc_zeroDim}; \pid{majorArc_successor}; \pid{majorArc_real};
\pid{modulatedAdjoint_identity}; \pid{structuredAdjoint_degreeLowering}.
\par}
The major-arc induction hypothesis must quantify over budgets; specializing its budget
before the induction would prevent its use on the enlarged supports.


\subsection{Remove the higher inputs, then prove lowest-input energy}
\label{patch:lowest-inverse}
For this subsection write a modulated correlation without inner-product notation as
\[
 \Gamma_l(g_0;f_1,\ldots,f_l;\zeta_{>l})
 =\int g_0(x)\mathbb E_t\prod_{i\leq l}f_i(x-P_i(t)e_i)
                   \ee{\sum_{i>l}\zeta_iP_i(t)}\,dx.
\]
The output factor is $g_0$, not its conjugate.  This convention eliminates implicit
conjugations in the identities below; the original four-linear form corresponds to
$g_0=\overline{f_0}$.

\begin{lemma}[Backward elimination of higher inputs]\label{patch:remove-high-inputs}
For an increasing degree list in $\{1,2,3\}$, a power-large unmodulated correlation
$\Gamma_k(g_0;f_1,\ldots,f_k)$ produces a fixed frequency vector $\zeta_{>1}$ and a
measurable $g'_0$ with $|g'_0|\leq|g_0|$ such that
\[
 |\Gamma_1(g'_0;f_1;\zeta_{>1})|\geq\delta^{O(1)}N^D.
\]
All inputs are $1$-bounded with budgeted support and admissible polynomials.
\end{lemma}
\begin{proof}
Induct backwards on $l=k,k-1,\ldots,2$ while retaining the original $f_1,\ldots,f_l$ and
one fixed vector $\zeta_{>l}$.  The base is the assumed correlation.
Suppose $\Gamma_l$ is power-large.  Define
\[
 F_l(y)=\mathbb E_t\overline{g_{0,l}(y+P_l(t)e_l)}
          \prod_{i<l}\overline{f_i(y+P_l(t)e_l-P_i(t)e_i)}
                           \ee{-\sum_{i>l}\zeta_iP_i(t)}.
\]
Then $\Gamma_l=\langle f_l,F_l\rangle$.  As in \eqref{patch:ma-adjoint}, replacing $f_l$
by $F_l$ gives a power-large correlation equal to $\|F_l\|_2^2$.
Apply \cref{patch:highest-control}, followed by \cref{patch:structured-degree} a bounded
number of times, to obtain power-large directional $U^2$ of $F_l$ at a radius comparable
to $N^{d_l}$ up to powers.

\paragraph{Measurable section frequencies.}
Let $y=x_{<l}$, $z=x_l$, and $w=x_{>l}$.
By \cref{patch:u2-fourier-selection} and popularity there is a measurable set
$Z\subseteq\{(y,w)\}$ of measure at least $\delta^{O(1)}N^{D-d_l}$ and a measurable
frequency $\lambda(y,w)$ for which
\begin{equation}\label{patch:high-section-coeff}
 \left|\int F_l(y,z,w)\ee{\lambda(y,w)z}\,dz\right|
                                  \geq\delta^{O(1)}N^{d_l}\qquad((y,w)\in Z).
\end{equation}
Popularity in $w$ gives a measurable set $W$ of measure at least
$\delta^{O(1)}N^{D-D_l}$ such that each $Z_w$ has measure at least
$\delta^{O(1)}N^{D_{l-1}}$.
Extend $\lambda$ off $Z$ by a fixed real number $\lambda_{\rm dummy}$ whose modulus is
larger than the major-arc bound used in the next paragraph.  This is permitted because that
bound is independent of the input phase values.  No Fourier-coefficient assertion is made
for the extended values off $Z$.

\paragraph{Use $\mathsf{MA}(l,l)$ to keep the selected frequencies small.}
For each $w\in W$, the integral over $y$ of the absolute coefficient with this extended
frequency is power-large.  Dualize by a measurable phase in $y$, expand $F_l$, and translate
by $P_l(t)e_l$.  This gives a power-large integral in $(y,z)$ with the $l-1$ active inputs
$\overline{f_i}$, the variable coefficient $-\lambda(y,w)$ of $P_l(t)$, and fixed coefficients
$-\zeta_i$ for $i>l$.
Select a $z$ section carrying a power-large correlation and apply the already proved
$\mathsf{MA}(l,l)$.  Its output budget is the same for every $w$.
In particular, the canonical set of $y$ with
\[
 N^{d_l}|\lambda(y,w)|+\sum_{i>l}N^{d_i}|\zeta_i|\leq u_C(\delta)
\]
is power-large.  Choose $|\lambda_{\rm dummy}|>u_C(\delta)N^{-d_l}$.
Every point of this canonical set must lie in $Z_w$, so the original estimate
\eqref{patch:high-section-coeff} holds there.  The full set of these $(y,w)$ is jointly
measurable, is power-large, and has $|\lambda(y,w)|\leq u_C(\delta)N^{-d_l}$.
This dummy-value step prevents the major-arc theorem from returning unrelated points on
which no useful Fourier coefficient was obtained.

\paragraph{Make the frequency constant.}
Partition that bounded frequency interval into cells of width $\tau N^{-d_l}$.
Choose $\tau>0$ so that changing $\lambda$ by at most $\tau N^{-d_l}$ changes the coefficient
in \eqref{patch:high-section-coeff} by at most one quarter of its lower bound.
Indeed $F_l$ is $1$-bounded and its $z$ support has length and absolute location bounded
by $u_c(\delta)N^{d_l}$; integrate
$|\ee{\lambda z}-\ee{\lambda_0 z}|\leq2\pi|\lambda-\lambda_0||z|$.
The chosen $\tau$ is a power of $\delta$ before the cell selection.
There are power-many cells, so a common center $\lambda_0$ works on a measurable set of
$(y,w)$ of power-large measure.  The coefficient at $\lambda_0$ is still pointwise power-large
on that set.  Choose a measurable $c(y,w)$, zero off that set and of modulus one on it,
such that
\[
 \int c(y,w)\int F_l(y,z,w)\ee{\lambda_0z}\,dz\,dy\,dw
\]
is nonnegative and power-large.
Expand $F_l$, translate by $P_l(t)e_l$, and conjugate the entire scalar integral.
The result is exactly $\Gamma_{l-1}$, with
\[
 g_{0,l-1}(x)=g_{0,l}(x)\overline{c(x_{\{l\}^c})}\ee{-\lambda_0x_l},
 \qquad \zeta_l=\lambda_0,
\]
and with the earlier fixed $\zeta_{>l}$ unchanged.  Thus
$|g_{0,l-1}|\leq|g_{0,l}|$, the support does not enlarge, and the remaining lower inputs are
still the original $f_i$, not modified products.
There are at most two backward steps.  All budget choices are explicit applications of
\cref{patch:budgets}, so the iteration yields the stated conclusion.
\end{proof}
The inequality $|g'_0|\leq|g_0|$ is the form needed here.  The paper states a modulus-equality
version in Lemma 4.71; no such stronger equality is needed or asserted in this patch.

\begin{theorem}[Lowest-input positive Fourier energy]\label{patch:lowest-energy}
Fix the increasing degree list and an input budget.  There is a structural integer $C$ such
that a power-large correlation $\Gamma_k$ as above, for $N\geq u_C(\delta)$, implies
\[
 \left\langle f_1,P_{L}^{(1)}f_1\right\rangle
       =\int\eta(\xi_1/L)|\widehat f_1(\xi)|^2\,d\xi
       \geq\ell_C(\delta)N^D,
       \qquad L=u_C(\delta)N^{-d_1}.
\]
The energy is real and nonnegative.  The theorem holds for every increasing sublist of
$\{1,2,3\}$, including lists of length one whose sole degree is two or three.
\end{theorem}
\begin{proof}
For $k\geq2$, apply \cref{patch:remove-high-inputs}.  For $k=1$ begin with the original
correlation.  Set
\[
 T f_1(x)=\mathbb E_t f_1(x-P_1(t)e_1)\ee{\sum_{i>1}\zeta_iP_i(t)}.
\]
Cauchy--Schwarz with the supported output factor gives
$\|Tf_1\|_2^2\geq\delta^{O(1)}N^D$.
Plancherel in all spatial variables gives
\[
 \|Tf_1\|_2^2=\int |m(\xi_1)|^2|\widehat f_1(\xi)|^2\,d\xi,
 \qquad m(v)=\mathbb E_t\ee{-vP_1(t)+\sum_{i>1}\zeta_iP_i(t)}.
\]
We have $|m|\leq1$ and $\|f_1\|_2^2\leq u_c(\delta)N^D$ for a known budget.
Choose a power-small $\varepsilon>0$ so that
$\varepsilon^2\|f_1\|_2^2$ is at most one quarter of the lower bound for $\|Tf_1\|_2^2$.
By \cref{patch:triangular}, $|m(v)|\geq\varepsilon$ implies
$|v|\leq R=u_{C'}(\delta)N^{-d_1}$.
Thus the Fourier energy in $|\xi_1|\leq R$ is power-large.
Choose $L\geq4R$.  The cutoff equals one on this band and is nonnegative elsewhere, so
\[
 \langle f_1,P_L^{(1)}f_1\rangle\geq
          \int_{|\xi_1|\leq R}|\widehat f_1(\xi)|^2d\xi
                              \geq\delta^{O(1)}N^D.
\]
Choose a single integer $C$ dominating the threshold, radius coefficient and exponent,
and lower-bound coefficient and exponent, so that $L=u_C(\delta)N^{-d_1}$ still has its
plateau on the same band.  This proves the exact statement.
\end{proof}
This last step uses the source's modulated-linear-average and scalar-oscillation argument,
but full Plancherel avoids an unnecessary final sectionwise choice of Fourier data.

\subsection{General selected indices: quantitative freezing and slicing}
\label{patch:all-indices}
The next theorem is the independent substitute for the source's Theorem 4.30 needed by the
public monomial wrappers.  Its proof is given here rather than taken as an axiom.

\begin{theorem}[Every-input positive energy, independently proved]
\label{patch:energy-core}
Fix an increasing list $1\leq d_1<\cdots<d_k\leq3$ and an integer input budget $c\geq1$.
There exists an integer $C\geq c$ such that the following holds for $0<\delta\leq1$ and
$N\geq u_C(\delta)$.  Suppose that all $P_i$ are admissible with budget $c$, the functions
$g_0,f_1,\ldots,f_k$ are $1$-bounded and supported in the budget-$c$ box, and
\[
 |\Gamma_k(g_0;f_1,\ldots,f_k)|\geq\ell_c(\delta)N^D.
\]
Then, for every $j\in\{1,\ldots,k\}$,
\begin{equation}\label{patch:all-energy}
 \left\langle f_j,P_{L_j}^{(j)}f_j\right\rangle
      \geq\ell_C(\delta)N^D,
       \qquad L_j=u_C(\delta)N^{-d_j}.
\end{equation}
\end{theorem}
\begin{proof}
Induct on $k$, with the statement quantified over all increasing degree lists and all input
budgets.  For $k=1$, or for $j=1$ at any $k$, use \cref{patch:lowest-energy}.
Suppose $k\geq2$ and $j>1$.

\paragraph{1. Replace the first input, and smooth that adjoint.}
Define its exact adjoint
\[
 F_1(y)=\mathbb E_t\overline{g_0(y+P_1(t)e_1)}
                  \prod_{i>1}\overline{f_i(y+P_1(t)e_1-P_i(t)e_i)}.
\]
The original scalar is $\langle f_1,F_1\rangle$.  Its size gives
$\|F_1\|_2^2\geq\delta^{O(1)}N^D$, and replacing $f_1$ by $F_1$ produces this positive
correlation.  All its inputs, including $F_1$, are $1$-bounded and supported in a budgeted
enlargement.  Apply \cref{patch:lowest-energy} to this new correlation to obtain
\[
 \langle F_1,P_{L_1}^{(1)}F_1\rangle\geq\kappa N^D,
 \qquad L_1=\delta^{-O(1)}N^{-d_1},\quad\kappa=\delta^{O(1)}.
\]
Set $G=P_{L_1}^{(1)}F_1$.  The adjoint identity gives the same lower bound for the
correlation with first input $G$ and other inputs the original $f_2,\ldots,f_k$.
The scalar is real by self-adjointness, but only its modulus will be used.

Let $M=\max(1,\|\check\eta\|_1)$ and $M_1=\| (\check\eta)'\|_1$.
The one-coordinate convolution representation gives
\[
 \|G\|_\infty\leq M,\qquad \|\partial_1G\|_\infty\leq M_1L_1.
\]
For bounded measurable $F_1$, differentiate the integrable smooth kernel; this justifies the
derivative bound and the one-variable Lipschitz estimate on every Borel section.
The output factor $g_0$ retains finite support even though $G$ itself need not have compact
support in coordinate one.

\paragraph{2. Freeze on one of finitely many equal intervals.}
Admissibility bounds $|P_1'(t)|\leq u_{c'}(\delta)N^{d_1-1}$ on $[0,N]$.
Partition $[0,N]$ into $K$ equal intervals $[a,a+n]$ with $n=N/K$.
On such an interval,
\[
 |G(x-P_1(t)e_1)-G(x-P_1(a)e_1)|
 \leq M_1L_1u_{c'}(\delta)N^{d_1-1}n
 \leq u_{c''}(\delta)/K.
\]
After integrating against $g_0$ and the other bounded factors, the total normalized
correlation error is at most $u_{c'''}(\delta)N^D/K$.
Choose
\[
 K=\max\{1,\lceil 4u_{c'''}(\delta)/\kappa\rceil\}.
\]
This is power-bounded above, and ensures error at most $\kappa N^D/4$.
The full average is the average of these $K$ normalized local averages.  Hence at least one
frozen interval has scalar modulus at least $\kappa N^D/2$.
Choose the first such interval in this finite family.  No discarded remainder interval is needed.

\paragraph{3. Expose the lower-dimensional polynomial family.}
For that $a$, put
\[
 Q_i(u)=P_i(a+u)-P_i(a),\qquad
 f_i^a(x)=f_i(x-P_i(a)e_i)\quad(i\geq2),
\]
\[
 g_0^a(x)=M^{-1}g_0(x)G(x-P_1(a)e_1).
\]
The frozen normalized correlation, divided by $M$, is exactly
\begin{equation}\label{patch:frozen-correlation}
 \int g_0^a(x)\mathbb E_{u\in[0,n]}
               \prod_{i=2}^k f_i^a(x-Q_i(u)e_i)\,dx.
\end{equation}
It is power-large in units $N^D$, and all its displayed functions are $1$-bounded.
The $Q_i$ have degrees $d_i$ and the same leading coefficients as $P_i$.
By the binomial formula and $0\leq a\leq N$, their degree-$r$ coefficients are bounded by
$u_{c'}(\delta)N^{d_i-r}$, hence by
$u_{c'}(\delta)K^{d_i-r}n^{d_i-r}$.
Since $K$ is power-bounded, these are admissible coefficients at scale $n$ with one new
structural budget.  The translated supports are contained in boxes with sides bounded by
$u_{c'}(\delta)N^{d_i}=u_{c'}(\delta)K^{d_i}n^{d_i}$ and therefore satisfy the same kind
of enlarged-support hypothesis at scale $n$.

\paragraph{4. Obtain a measurable set of large sections.}
For $v=x_1$, let $B(v)$ be the integral in the remaining coordinates in
\eqref{patch:frozen-correlation}.  It is measurable by Fubini, vanishes unless $v$ lies in
the original coordinate-one support interval, and
$|B(v)|\leq u_{c'}(\delta)N^{D-d_1}$.
Since $\int|B(v)|dv\geq|\int B(v)dv|\geq\delta^{O(1)}N^D$, popularity produces a
canonical measurable set $X_1$ of measure at least $\delta^{O(1)}N^{d_1}$ on which
\[
 |B(v)|\geq\delta^{O(1)}N^{D-d_1}\geq\delta^{O(1)}n^{D-d_1}.
\]
All hypotheses of the $(k-1)$-dimensional induction theorem now hold with the increasing
list $(d_2,\ldots,d_k)$ and one common input budget, independent of $v\in X_1$.
This includes lists starting with degree two or three, not only lists starting with one.
Because $n=N/K$ and $K\leq u_{c'}(\delta)$, choosing the original threshold $N$ larger
by the product of the two known budget bounds ensures the required lower threshold for $n$.

\paragraph{5. Apply the induction hypothesis and integrate positive energies.}
The induction hypothesis gives, for every $v\in X_1$, a lower bound
$\delta^{O(1)}n^{D-d_1}$ for the positive energy of the section of $f_j^a$ at the same
radius $L'=u_{C'}(\delta)n^{-d_j}$.  Translation by $P_j(a)e_j$ does not change its
coordinate-one section and preserves the directional Fourier energy in the remaining
coordinates.  Thus this is also the section energy of the original $f_j$.
Since $n=N/K$ and $K$ is power-bounded, the lower bound is
$\delta^{O(1)}N^{D-d_1}$.
Integrate over $X_1$ and use nonnegativity on its complement.  This gives
\[
 \langle f_j,P_{L'}^{(j)}f_j\rangle\geq\delta^{O(1)}N^D,
 \qquad L'\leq u_{C''}(\delta)N^{-d_j}.
\]
The radius is common to all sections; no varying-radius integral is being identified with
one projection.
Finally choose $L_j=u_C(\delta)N^{-d_j}\geq2L'$.  On the support of
$\eta(\xi_j/L')$, the larger cutoff equals one, so
$\eta(\xi_j/L_j)\geq\eta(\xi_j/L')$ pointwise.
The positive energy therefore only increases.  A common maximum of the finitely many
budgets handles all $j$, thresholds, and the degree-list induction.
This proves \eqref{patch:all-energy}.
\end{proof}

\paragraph{Why this is not the old circular proof.}
\Cref{patch:energy-core} uses the independently proved \cref{patch:lowest-energy}, the
adjoint identities, freezing, and lower-dimensional instances of itself.
It does not use the public \pid{lem:pet-reduction}, the old \pid{lem:degree-lowering-zero},
or the public \pid{thm:real-inverse}.  The only PET premise used in proving
\cref{patch:lowest-energy} was proved explicitly in \cref{patch:pet-producer}.
\paragraph{Lean nodes.}
\mbox{}\par
{\small\raggedright
\pid{higherInput_eliminate}; \pid{lowestInput_positiveEnergy};
\pid{projectedAdjoint_lipschitz}; \pid{freeze_parameter_interval};
\pid{shiftedPolynomial_admissible}; \pid{largeSection_set};
\pid{positiveEnergy_slice}; \pid{real_everyInput_positiveEnergy}.
\par}


\section{The public every-input conclusion and its actual cube}
\label{patch:public-completion}
We now return to exactly the monomial system and normalizer of the original manuscript.
Fix a support parameter $C_0\geq1$, enlarge it to an integer if necessary, and put
\[
 B_N=I_N(C_0)=[-C_0N,C_0N]\times[-C_0N^2,C_0N^2]\times[-C_0N^3,C_0N^3],
 \qquad\beta=(2C_0)^3,
\]
so that $|B_N|=\beta N^6$.

\begin{corollary}[Independent monomial projected energy]\label{patch:monomial-energy}
There are integers $a,b\geq1$, depending only on $C_0$ and the fixed cutoff, such that if
$0<\delta\leq1$, $N\geq a\delta^{-a}$, all four $f_i$ are $1$-bounded and supported in
$B_N$, and $|\Lambda_N(f_0,f_1,f_2,f_3)|\geq\delta N^6$, then, for every $j\in\{1,2,3\}$,
\begin{equation}\label{patch:projected-energy}
 E_{L_j}(f_j):=N^{-6}\|P_{L_j}^{(j)}f_j\|_2^2
       \geq a^{-1}\delta^a,
       \qquad L_j=b\delta^{-b}N^{-j}.
\end{equation}
\end{corollary}
\begin{proof}
Use \cref{patch:energy-core} with $P_i(t)=-t^i$, $g_0=\overline{f_0}$, and a fixed
input budget $c\geq\max(1,\lceil C_0\rceil)$.
The input support and admissibility conditions hold for every $0<\delta\leq1$, and
$\delta\geq\ell_c(\delta)$.  The minus sign matches the source-style translation
$x-P_i(t)e_i=x+t^ie_i$ exactly.
Let $C$ be the returned output budget.  The positive energy at
$L_0=C\delta^{-C}N^{-j}$ is at least $C^{-1}\delta^C N^6$.
Set $a=C$ and $b=2C$.  Then $L_j\geq2L_0$ and
\[
 |\eta(\xi_j/L_j)|^2\geq\eta(\xi_j/L_0)
\]
pointwise: wherever the right side is nonzero the left side equals one.
Plancherel yields \eqref{patch:projected-energy}.  No source theorem is a new hypothesis
here; \cref{patch:energy-core} was proved in the preceding section.
\end{proof}

\begin{lemma}[Low-frequency energy constructs the order-two cube]
\label{patch:energy-to-u2}
Let $f$ be $1$-bounded and supported in $B_N$, let $L>0$ satisfy $LN^j\geq1$, and put
$H=1/(8L)$.  With the \emph{existing} local-uniformity definition,
\begin{equation}\label{patch:energy-u2-quantitative}
 \mathcal U^2_{j,H}(f)^4
   =\operatorname{locUnifPow}(f;j,H,2)
   \geq\frac{E_L(f)^2}{32\beta}.
\end{equation}
\end{lemma}
\begin{proof}
Write $Q_s=Q_{s,H,N^6}^j(f)$, so $Q_s$ is the public, not the support-volume-normalized,
power.  Define $T_Hu(x)=H^{-1}\int_0^H u(x+te_j)dt$.
By \cref{patch:fejer-squares},
\begin{equation}\label{patch:public-u1}
 Q_1=N^{-6}\|T_Hf\|_2^2.
\end{equation}
The multiplier of $T_H$ is $m_H(v)=H^{-1}\int_0^H\ee{vt}dt$.
For $|v|\leq L/2$, the elementary exponential estimate gives
\[
 |m_H(v)-1|\leq H^{-1}\int_0^H2\pi|v|t\,dt
                  =\pi H|v|\leq\pi/16<1/4.
\]
In particular $|m_H(v)|\geq1/2$ there.  Outside that band $\eta(v/L)=0$.
Thus $|\eta(v/L)|^2\leq4|m_H(v)|^2$ for every $v$; Plancherel and
\eqref{patch:public-u1} imply
\begin{equation}\label{patch:public-energy-u1}
 Q_1\geq E_L(f)/4.
\end{equation}

For each real $z$ put
\[
 g_z(x)=f(x)\overline{f(x+ze_j)},\qquad c(z)=\int g_z(x)dx.
\]
The zero vertex $f(x)$ supports $g_z$ in $B_N$.
If $S=2C_0N^j$ is its $j$-th side length, $T_Hg_z$ is supported in a box of volume
$V_H=\beta N^6(1+H/S)$.  Its integral equals $c(z)$, so
\begin{equation}\label{patch:public-support-cs}
 |c(z)|^2\leq V_H\|T_Hg_z\|_2^2.
\end{equation}
Expanding this square and the definition of $g_z$ gives the literal identity
\begin{align}
 \|T_Hg_z\|_2^2
 &=\int\kappa_H(w)\int g_z(x)\overline{g_z(x+we_j)}\,dx\,dw\notag\\
 &=\int\kappa_H(w)\int
 f(x)\overline{f(x+ze_j)}\overline{f(x+we_j)}
                  f(x+(z+w)e_j)\,dx\,dw.\label{patch:public-literal-cube}
\end{align}
Consequently
\begin{equation}\label{patch:public-u2-squares}
 Q_2=N^{-6}\int\kappa_H(z)\|T_Hg_z\|_2^2\,dz.
\end{equation}
Jensen for the probability measure $\kappa_H(z)dz$, followed by
\eqref{patch:public-support-cs}, now yields
\[
 Q_1^2=\left|N^{-6}\int\kappa_H(z)c(z)dz\right|^2
 \leq N^{-12}\int\kappa_H(z)|c(z)|^2dz
 \leq (V_H/N^6)Q_2.
\]
Since $LN^j\geq1$, $H\leq N^j/8\leq S$ and $V_H/N^6\leq2\beta$.
Combine this with \eqref{patch:public-energy-u1} to obtain
$Q_2\geq E_L(f)^2/(32\beta)$.
All integrals are absolutely integrable by boundedness, finite support, and compactly supported
shift measures.  In particular, the argument has not identified a $U^2$ power with the
first-order Fourier energy; the intervening support Cauchy--Schwarz step is indispensable.
\end{proof}

\paragraph{Exact data supplied to \texttt{cfgHeadInt\_cube}.}
Use order two and the four atoms indexed by $\omega\in\{0,1\}^2$:
\[
 \begin{array}{c|c|c|c}
 \omega &\text{origin function}&\text{translation}&\text{conjugation flag}\\\hline
 (0,0)&f&0&0\\
 (1,0)&f&ze_j&1\\
 (0,1)&f&we_j&1\\
 (1,1)&f&(z+w)e_j&0
 \end{array}
\]
The base variable is $x\in\mathbb R^3$, its normalizer is $N^{-6}$, both increment measures
are independently $\kappa_H(t)dt$, and the amplitude is exactly one.
Equation \eqref{patch:public-literal-cube} proves equality with these four atoms;
\eqref{patch:public-u2-squares} integrates the remaining Fej\'er coordinate.
Only after this equality is proved does one invoke the existing \pid{cfgHeadInt_cube}.
There is no residual $t$, no additional selected-input certificate, and no product hidden
in the origin function.  The bound $32\beta$ stays outside the configuration.

\begin{lemma}[Every-input local uniformity for the monomial average]
\label{lem:pet-reduction}
For every fixed $C_0\geq1$, there are $s_*,C_*\in\mathbb N$ with $s_*=2$, $C_*\geq2$,
and $c_*>0$, such that the following holds.  If $0<\delta\leq1$, $N\geq C_*\delta^{-C_*}$,
all four $f_i$ are $1$-bounded and supported in $I_N(C_0)$, and
$|\Lambda_N(f_0,f_1,f_2,f_3)|\geq\delta N^6$, then for every $j\in\{1,2,3\}$ there exists
\[
 c_*\delta^{C_*}N^j\leq H_j\leq\delta^{-C_*}N^j
\]
with $\mathcal U^{s_*}_{j,H_j}(f_j)\geq c_*\delta^{C_*}$.
\end{lemma}
\begin{proof}
Let $a,b$ be the integers from \cref{patch:monomial-energy} and choose explicitly
\[
 H_j=\frac{\delta^b}{8b}N^j,
 \quad C_*\geq\max(a,b,2),
 \quad0<c_*\leq\min\left\{\frac1{8b},(32\beta a^2)^{-1/4},1\right\}.
\]
The scale threshold implies $N\geq a\delta^{-a}$, so
\cref{patch:monomial-energy,patch:energy-to-u2} give
\[
 \operatorname{locUnifPow}(f_j;j,H_j,2)
       \geq\frac{\delta^{2a}}{32\beta a^2},\qquad
 \mathcal U^2_{j,H_j}(f_j)
       \geq(32\beta a^2)^{-1/4}\delta^{a/2}
       \geq c_*\delta^{C_*}.
\]
Here $L_jN^j=b\delta^{-b}\geq1$ verifies the scale hypothesis in the bridge.
Since $0<\delta\leq1$ and $C_*\geq b$, the displayed radius satisfies both requested
bounds.  There is one common $C_*,c_*$ for all three inputs.
In Lean prove the fourth-power inequality first and apply nonnegative-root monotonicity
only once, at the end.
\end{proof}
The old label is retained for reference compatibility, but the proof no longer says that an
unspecified every-input PET recursion terminates in the desired cube.  Its every-input
conclusion is the conclusion actually needed by the manuscript, now with an explicit order.

\begin{theorem}[Specialized real inverse theorem]\label{thm:real-inverse}
For each $C_0\geq1$ there is an integer $C\geq1$ such that if $0<\delta\leq1$,
$N\geq C\delta^{-C}$, the four functions are $1$-bounded and supported in $I_N(C_0)$,
and $|\Lambda_N(f_0,f_1,f_2,f_3)|\geq\delta N^6$, then for every $j$ there is a
$1$-bounded function $h_j$, supported in $I_N(C_0)$, for which
\[
 |\langle f_j,P_L^{(j)}h_j\rangle|\geq C^{-1}\delta^C N^6,
                   \qquad L=C\delta^{-C}N^{-j}.
\]
One may take $h_j=f_j$.
\end{theorem}
\begin{proof}
Use \cref{patch:monomial-energy}, not \cref{lem:pet-reduction}.
Let $L_0=b\delta^{-b}N^{-j}$ and choose an integer $C\geq\max(a,2b,1)$.
Then $L\geq2L_0$ and the cutoff plateau gives
$\eta(\xi_j/L)\geq|\eta(\xi_j/L_0)|^2$ for every $\xi$.
Plancherel and positivity imply
\[
 \langle f_j,P_L^{(j)}f_j\rangle
       \geq\|P_{L_0}^{(j)}f_j\|_2^2
       \geq a^{-1}\delta^aN^6
       \geq C^{-1}\delta^CN^6.
\]
This real nonnegative scalar has the same modulus.  The witness $h_j=f_j$ already has
exactly the required boundedness and support.  There is no kernel-tail or phase-selection
step and no use of the old degree-lowering lemma.
\end{proof}

\paragraph{The deleted degree-lowering assertion.}
The original \pid{lem:degree-lowering-zero} and its proof are not part of the replacement.
Its proposed argument applied degree lowering to a generic selected input and identified a
$U^2$ expression with the wrong Fourier energy.  Here degree lowering is used only for the
structured adjoints in \cref{patch:conditional-degree,patch:structured-degree}; positive
low-frequency energy is proved subsequently by the scalar oscillatory multiplier argument.
Both public conclusions are downstream of that same independently proved energy theorem.

\paragraph{Lean nodes for the public wrappers.}
\mbox{}\par
{\small\raggedright
\pid{monomial_positiveEnergy}; \pid{projection_energy_le_four_u1};
\pid{u1_sq_le_support_mul_u2}; \pid{energy_to_locUnifPow_two};
\pid{everyInput_localUniformity}; \pid{realInverse_from_positiveEnergy}.
\par}
The hypotheses of \pid{everyInput_localUniformity} must be exactly those in the public
lemma, with no \pid{SourceInverse}, \pid{MajorArcProperty}, or terminal-state argument.


\section{Implementation order and application to the original file}
\label{patch:implementation}
\subsection{Acyclic implementation order}
The following table is a topological order of the retained proof arguments.  A row refers to
the mathematical results above, not to a proposed axiom with that name.
\begin{center}\small
\begin{tabular}{>{\raggedright\arraybackslash}p{.06\linewidth}>{\raggedright\arraybackslash}p{.47\linewidth}>{\raggedright\arraybackslash}p{.35\linewidth}}
\hline
Stage & Declarations to implement & Earlier mathematical dependencies \\\hline
1 & Budget arithmetic; measurability; Fej\'er square identity and cube recursion & Background measure theory and finite products \\
2 & Signed integrated VdC; cylinder-CS state; scalar oscillation and triangular coefficient bound & Stage 1; scalar calculus \\
3 & Protected leading invariant; specified PET pivot; uniform termination; affine block cube; sublevel exclusion and uniformization & Stages 1 and 2 \\
4 & $U^2$/Fourier selection; dual--difference interchange; missing-coordinate phase removal; dummy phase & Stages 1 and 2; Plancherel \\
5 & Conditional structured degree lowering, as an implication from $\mathsf{MA}(m,l)$ & Stage 4 \\
6 & Base and successor major-arc induction; unconditional structured-adjoint lowering & Stages 2--5; smaller $m$ in the induction \\
7 & Backward elimination of higher inputs; lowest-input positive energy & Stages 3, 4 and 6 \\
8 & Projected-adjoint Lipschitz bound; short-interval freezing; lower-dimensional slicing; every-input positive energy & Stage 7 and smaller ambient dimension \\
9 & Monomial energy specialization; concrete two-dimensional cube; every-input local uniformity; public inverse theorem & Stages 1 and 8 \\\hline
\end{tabular}
\end{center}
There are two distinct inductions in Stages 6--8: first the number $m-1$ of functions in a
major-arc correlation, and only later the ambient dimension in the all-input energy theorem.
They are not interchangeable.  The backward elimination in Stage 7 is a separate finite
induction on the number of remaining active inputs.

\subsection{Representation details that matter in Lean}
The polynomial state must retain the following data: the ordered normal polynomial list;
a list of origin-labelled, translated/conjugated factors for each block; the current spatial
block; a marked singleton protected block; its original degree and leading coefficient;
the formal nonzero integer polynomials witnessing \eqref{patch:lead-invariant}; and numerical
length, coefficient, support and loss budgets.  A Boolean conjugation flag is enough.
For the ambient space one may use the project's existing Euclidean/product-measure model.
No particular Mathlib constructor or unprovided project field name is presumed.
Use the proof-local $Q_{s,H,V}^j$ in each ambient dimension $1,2,3$ (and the explicitly
trivial point measure for dimension zero).  The public dimension-three specialization
has normalizer $N^6$ and is identified with the existing \pid{locUnifPow}.
At a polynomial update embed the old shift-variable ring into the new ring and introduce
one genuinely fresh final variable.  Coefficient extraction in that variable is what proves
noncancellation in \cref{patch:pet-update}.

The essential equalities are \eqref{patch:cube-succ},
\eqref{patch:affine-cube-invariant},
\eqref{patch:affine-terminal}, the four-step phase identity, the cylinder states in
\cref{patch:cylinder}, the exact adjoints \eqref{patch:ma-adjoint} and
\eqref{patch:ma-smaller-pattern}, the frozen identity \eqref{patch:frozen-correlation},
and the literal four-factor identity \eqref{patch:public-literal-cube}.
Each must be established as an equality before invoking an analytic inequality.

Use lists or finite indexed products with multiplicity for factors.  A finite set of
polynomials is suitable for counting leading-coefficient classes, but is not a replacement
for the factor list.  Grouping equal polynomial arguments multiplies their attached functions;
it does not discard an atom.  The invariant proves that no such grouping alters the protected
singleton before the terminal stage.

Use explicit real integrals of nonnegative quantities when applying monotonicity, Jensen,
or a kernel bound.  Complex cube integrals first become square averages by
\cref{patch:fejer-squares}.  The support normalizer $V$, section normalizer $S$, and public
normalizer $N^6$ should be separate arguments, with separate conversion lemmas.
For a positive Fourier-energy pairing, store its real part in Lean and prove that its
imaginary part is zero before using real inequalities.  Prove fourth-power bounds before taking fourth roots.  Treat zero norms,
zero polynomial oscillation coefficients, and null sections in their stated separate branches.

\subsection{Scoped replacement map}
Apply these instructions to the \emph{current} manuscript, preserving all previously completed
repairs outside this scope.  Use unique theorem and subsection labels rather than shifted
line numbers.  This deliverable is a standalone blueprint, not a replacement for the entire
manuscript; it does not include the earlier combined diff or application script.

\paragraph{Primary edits: original lines 2183 and 2193.}
Inside \pid{lem:fejer-vdc}, retain the scalar inequality and its existing Lean proof
\pid{fejer_vdc'}.  Delete only the unrestricted last assertion and its unsupported final
proof paragraph.  The scalar statement reproduced in this document recalls that interface;
it does not ask for a second formalization of the already proved inequality.
Separate the new signed integrated removal lemma, coordinate-removal induction, and
polynomial arguments into their own precisely scoped declarations.

\paragraph{Retain the existing definitions and bridge.}
Keep \pid{def:local-uniformity}, \pid{locUnifPow}, and \pid{cfgHeadInt_cube}.
Their definitions are recalled here so the blueprint can be read independently.
When integrating the proof, reuse these declarations rather than duplicating them.
The proof-local normalizer $V$ and lower-dimensional cube quantities are auxiliary;
the public quantity still has normalizer $N^6$.
Likewise, retain a single implementation of the degree-at-most-three scalar oscillation
result used here, rather than introducing duplicate theorem names.

\paragraph{Necessary dependent edits: the PET and inverse proof path.}
Insert the conventions and independent inverse-module arguments before their public
wrappers, in the topological order above.  Within the existing
\pid{sec:internal-kosz-adjoint}, replace the proof and title of \pid{lem:pet-reduction}
by the every-input local-uniformity result in \cref{lem:pet-reduction}.
Its proof uses the positive-energy theorem and explicit order-two cube, not an assumed
every-input terminal configuration.  Replace the proof of \pid{thm:real-inverse}
by the separate positive-energy wrapper in \cref{thm:real-inverse}.

Remove the old \pid{lem:degree-lowering-zero} from this replaced proof path.
The structured degree-lowering declarations proved here have explicit adjoint hypotheses;
they are not substitutes for a generic degree-lowering rule on arbitrary selected inputs.
Check any remaining callers in the actual Lean project against those hypotheses rather
than silently reusing the old assertion.

The existing labels \pid{sec:internal-kosz-adjoint}, \pid{def:local-uniformity},
\pid{lem:fejer-vdc}, \pid{lem:real-polynomial-oscillation}, \pid{lem:pet-reduction},
and \pid{thm:real-inverse} can be retained, once each, for reference compatibility.
Retaining a LaTeX label does not retain the old unrestricted mathematical claim.
The Hahn--Banach subsection and later callers continue to use the public inverse conclusion;
they are not rewritten in this scope revision.

\paragraph{Implementation map, sources, and protected earlier work.}
Update only the inverse-theorem portion of the manuscript's closing dependency map to use
the order above.  Add the relevant source locators from Sections 4.2--4.5 of the supplied
paper without removing bibliography information needed elsewhere.
Leave the completed subunit repair, its public wrapper, its constants, and its dependency-map
entries unchanged.  No part of the earlier stopping-time/subunit proof is replaced by this
document, and no other manuscript assertions are audited here.

\subsection{What counts as fully discharged in the Lean project}
The exported theorem types must have the original monomial input hypotheses, with no extra
certificate asserting the hard conclusion.  A theorem of type
$\texttt{SourceInverse}\to\texttt{EveryInputUniformity}$ does not discharge the public result.

After implementation, build the Lean project, inspect the exported theorem types, and audit
their transitive axiom dependencies:
\begin{verbatim}
#print everyInput_localUniformity
#print axioms everyInput_localUniformity
#print axioms realInverse_from_positiveEnergy
\end{verbatim}
These are proposed exported names; use the project's actual names.
No \pid{sorryAx} or additional analytic axiom should occur.  Report Lean's standard logical
axioms separately.  The official manual explains the audit command \cite{LeanValidation}.
Typesetting and reference checks are not a Lean compilation; none is claimed here.

\begin{thebibliography}{9}
\bibitem{KoszEtAl}
D.~Kosz, M.~Mirek, S.~Peluse, R.~Wan, and J.~Wright,
\emph{The multilinear circle method and a question of Bergelson},
supplied arXiv:2411.09478v4 PDF (105 pages).
The source organization used here is Sections 4.2--4.5.
The mathematical adaptations and additional derivations are identified in the text.
\bibitem{LeanValidation}
Lean Language Reference, \emph{Validating a Lean Proof}, official online manual,
accessed 15 September 2026.  The \texttt{\#print axioms} command and its interpretation.
\end{thebibliography}
\end{document}
